\documentclass[a4paper,12pt]{article}

\usepackage[a4paper,margin=2.5cm]{geometry}
\usepackage{iftex}
\ifPDFTeX
  \usepackage[T2A]{fontenc}
  \usepackage[utf8]{inputenc}
  \usepackage[russian]{babel}
\else
  \usepackage[russian]{babel}
  \usepackage{fontspec}
  \setmainfont{Times New Roman}
\fi

\usepackage{amsmath}
\usepackage{amssymb}
\usepackage{enumitem}

\setlength{\parindent}{0pt}
\setlength{\parskip}{0.45em}
\setlist[itemize]{leftmargin=1.6em,itemsep=0.25em,topsep=0.25em}
\setlist[enumerate]{leftmargin=1.8em,itemsep=0.25em,topsep=0.25em}

\newcommand{\R}{\mathbb{R}}
\newcommand{\Answer}{\textbf{Ответ.}\ }
\newcommand{\Solution}{\textbf{Решение.}\ }

\begin{document}

\begin{center}
{\LARGE Коллоквиум по курсу}\\[0.2em]
{\LARGE «Уравнения математической физики»}\\[0.8em]
{\large Коллоквиум 2. Разделы 1--2}
\end{center}

\vspace{0.8em}
\textit{Конспект составлен по вопросам из приложенного списка и по пособию Г.\,В.~Алексеева «Классические методы математической физики. Часть 1». В формулировках и обозначениях максимально сохранена логика пособия. Для совместимости с Overleaf мнемонические блоки оформлены текстовыми заголовками без emoji.}

\section*{1. Уравнения в частных производных первого порядка}

\subsection*{1(a). Одномерное уравнение переноса с постоянным коэффициентом. Характеристики}

\Answer
В пособии сначала рассматривается простейшее одномерное уравнение
\[
\frac{\partial u}{\partial t}+\frac{\partial u}{\partial x}=0,
\]
а затем его более общий вариант
\[
\frac{\partial u}{\partial t}+a\frac{\partial u}{\partial x}=0,
\qquad a=\mathrm{const}>0.
\]

\textbf{Определение.} Уравнение
\[
u_t+a u_x=0
\]
называется \textbf{одномерным уравнением переноса} с постоянным коэффициентом. Здесь:
\begin{itemize}
  \item $u=u(x,t)$ --- искомая функция;
  \item $x$ --- пространственная переменная;
  \item $t$ --- время;
  \item $a$ --- постоянная скорость переноса.
\end{itemize}

\textbf{Физический смысл.} Если $a>0$, то профиль $u$ переносится вправо со скоростью $a$; если $a<0$, то влево со скоростью $|a|$.

\textbf{Характеристики.} Характеристиками называют такие кривые на плоскости $(x,t)$, вдоль которых уравнение в частных производных сводится к обыкновенному дифференциальному уравнению.

Для уравнения
\[
u_t+a u_x=0
\]
характеристики ищутся в виде кривых $x=x(t)$, вдоль которых полная производная функции $u(x(t),t)$ принимает простой вид. Для этого требуем
\[
\frac{dx}{dt}=a.
\]

\textbf{Пошаговое нахождение характеристик:}
\begin{enumerate}
  \item Записываем характеристическое уравнение:
  \[
  \frac{dx}{dt}=a.
  \]
  \item Поскольку $a=\mathrm{const}$, интегрируем:
  \[
  x=at+C.
  \]
  \item Переносим $at$ влево:
  \[
  x-at=C.
  \]
\end{enumerate}

Следовательно, \textbf{характеристики --- это семейство прямых}
\[
x-at=C.
\]

В частном случае $a=1$ получаем прямые
\[
x-t=C.
\]

\textbf{Геометрический смысл.} Каждая характеристика --- это траектория, вдоль которой переносится значение решения. Номер характеристики задается постоянной $C$.

\paragraph{Мнемоническое правило / Суть на пальцах.}
Уравнение переноса ничего не «рождает» и не «размывает»: оно просто сдвигает профиль вдоль прямых $x-at=\mathrm{const}$. Характеристика --- это дорожка, по которой значение $u$ едет без изменения.

\subsection*{1(b). Общее решение однородного уравнения переноса}

\Answer
Рассмотрим однородное уравнение
\[
u_t+a u_x=0,
\qquad a=\mathrm{const}.
\]

\textbf{Идея метода.} Нужно исследовать поведение решения вдоль характеристики
\[
x=x(t)=at+C.
\]

\textbf{Шаг 1.} Рассмотрим функцию
\[
t\longmapsto u(x(t),t).
\]
Это сложная функция одной переменной $t$.

\textbf{Шаг 2.} Вычислим ее полную производную по правилу дифференцирования сложной функции:
\[
\frac{d}{dt}u(x(t),t)=u_t(x(t),t)+u_x(x(t),t)\frac{dx}{dt}.
\]

\textbf{Шаг 3.} Так как вдоль характеристики
\[
\frac{dx}{dt}=a,
\]
то
\[
\frac{d}{dt}u(x(t),t)=u_t+a u_x.
\]

\textbf{Шаг 4.} Но по самому уравнению
\[
u_t+a u_x=0.
\]
Значит,
\[
\frac{d}{dt}u(x(t),t)=0.
\]

\textbf{Шаг 5.} Если производная вдоль характеристики равна нулю, то функция постоянна на этой характеристике:
\[
u(x(t),t)=\mathrm{const}.
\]

\textbf{Шаг 6.} Поскольку характеристика задается уравнением
\[
x-at=C,
\]
то значение решения зависит только от комбинации $x-at$. Следовательно,
\[
u(x,t)=\Phi(x-at),
\]
где $\Phi$ --- произвольная дифференцируемая функция одной переменной.

\textbf{Проверка подстановкой.}

Вычислим производные:
\[
u_x=\Phi'(x-at),
\qquad
u_t=-a\Phi'(x-at).
\]

Тогда
\[
u_t+a u_x=-a\Phi'(x-at)+a\Phi'(x-at)=0.
\]

Следовательно, \textbf{общее решение однородного уравнения переноса} имеет вид
\[
u(x,t)=\Phi(x-at).
\]

Для частного уравнения
\[
u_t+u_x=0
\]
получаем
\[
u(x,t)=\Phi(x-t).
\]

\paragraph{Мнемоническое правило / Суть на пальцах.}
Общее решение всегда имеет вид «начальный профиль, сдвинутый на $at$». Если нет правой части и нет диффузии, то форма графика не меняется вообще: он только едет.

\subsection*{1(c). Задача Коши для уравнения переноса на всей прямой}

\Answer
Рассмотрим задачу Коши
\[
u_t+a u_x=0,
\qquad (x,t)\in \R\times(0,+\infty),
\]
\[
u(x,0)=\varphi(x),
\qquad x\in\R,
\]
где, например, $\varphi\in C^1(\R)$.

\textbf{Что требуется сделать.} Нужно из общего решения
\[
u(x,t)=\Phi(x-at)
\]
выбрать ту функцию $\Phi$, которая удовлетворяет начальному условию.

\textbf{Шаг 1.} Подставляем $t=0$:
\[
u(x,0)=\Phi(x).
\]

\textbf{Шаг 2.} Сравниваем с начальным условием:
\[
\Phi(x)=\varphi(x).
\]

\textbf{Шаг 3.} Следовательно,
\[
\Phi=\varphi.
\]

\textbf{Итак, решение задачи Коши имеет вид}
\[
u(x,t)=\varphi(x-at).
\]

\textbf{Почему это решение единственно.} Через каждую точку $(x,t)$ проходит ровно одна характеристика
\[
x-at=C.
\]
Она пересекает начальную прямую $t=0$ ровно в одной точке
\[
(\xi,0),
\qquad \xi=x-at.
\]
Значение решения в точке $(x,t)$ обязано совпадать со значением начальной функции в точке пересечения:
\[
u(x,t)=\varphi(\xi)=\varphi(x-at).
\]
Значит, никакой свободы выбора уже не остается.

\textbf{Замечание о неоднородном уравнении.} Для задачи
\[
u_t+a u_x=f(x,t),
\qquad
u(x,0)=\varphi(x),
\]
метод характеристик дает формулу
\[
u(x,t)=\varphi(x-at)+\int_0^t f\bigl(x-a(t-\tau),\tau\bigr)\,d\tau.
\]
Она получается интегрированием равенства
\[
\frac{d}{dt}u(x(t),t)=f(x(t),t)
\]
вдоль характеристики $x(t)=x_0+at$.

\paragraph{Мнемоническое правило / Суть на пальцах.}
Чтобы узнать $u(x,t)$, надо от точки $(x,t)$ пройти назад по характеристике до начальной прямой. Куда пришли, то начальное значение и переносится в момент $t$.

\subsection*{1(d). Начально-краевая задача на полупрямой. Метод характеристик}

\Answer
В пособии подробно разобрана начально-краевая задача в полосе
\[
Q=(0,l]\times(0,T]
\]
для уравнения
\[
u_t+a u_x=f(x,t),
\qquad a=\mathrm{const}>0,
\]
с условиями
\[
u(x,0)=\varphi(x),
\qquad x\in[0,l],
\]
\[
u(0,t)=g(t),
\qquad t\in[0,T].
\]

Для \textbf{полупрямой} $x>0$ постановка полностью аналогична:
\[
u_t+a u_x=f(x,t),
\qquad x>0,\quad t>0,
\]
\[
u(x,0)=\varphi(x),
\qquad x\ge 0,
\]
\[
u(0,t)=g(t),
\qquad t\ge 0.
\]

\textbf{Почему при $a>0$ нужно задавать граничное условие именно в точке $x=0$.}
Характеристики имеют вид
\[
x-at=C.
\]
При $a>0$ они направлены вправо, поэтому часть характеристик входит в область $x>0$ через границу $x=0$. Чтобы задать решение на этих входящих характеристиках, нужно знать данные на границе.

\textbf{Алгоритм метода характеристик.}

\textbf{Шаг 1.} Берем произвольную точку $(x,t)$ в полуплоскости $x>0$, $t>0$.

\textbf{Шаг 2.} Проводим через нее характеристику
\[
x(\tau)=x-a(t-\tau).
\]
Это та же самая прямая $x-a\tau=\mathrm{const}$, записанная так, чтобы при $\tau=t$ мы попадали в точку $(x,t)$.

\textbf{Шаг 3.} Смотрим, куда эта характеристика приходит при движении назад по времени.

Возможны два случая.

\textbf{Случай 1:} характеристика пересекает начальную прямую $t=0$ раньше, чем границу $x=0$.

Это происходит, когда
\[
x-at\ge 0.
\]
Тогда точка пересечения с начальной прямой имеет вид
\[
(x-at,0),
\]
и вдоль характеристики получаем
\[
u(x,t)=\varphi(x-at)+\int_0^t f\bigl(x-a(t-\tau),\tau\bigr)\,d\tau.
\]

\textbf{Случай 2:} характеристика раньше попадает на границу $x=0$.

Это происходит, когда
\[
x-at<0.
\]
Тогда момент пересечения с границей находится из уравнения
\[
0=x-a(t-\tau_0),
\]
откуда
\[
\tau_0=t-\frac{x}{a}.
\]
Следовательно,
\[
u(x,t)=g\!\left(t-\frac{x}{a}\right)+\int_{t-x/a}^t f\bigl(x-a(t-\tau),\tau\bigr)\,d\tau.
\]

\textbf{Итоговая формула на полупрямой при $a>0$:}
\[
u(x,t)=
\begin{cases}
\displaystyle
\varphi(x-at)+\int_0^t f\bigl(x-a(t-\tau),\tau\bigr)\,d\tau,
& x\ge at,
\\[1.2em]
\displaystyle
g\!\left(t-\frac{x}{a}\right)+\int_{t-x/a}^t f\bigl(x-a(t-\tau),\tau\bigr)\,d\tau,
& 0<x<at.
\end{cases}
\]

Для однородного уравнения $f\equiv 0$ формула упрощается:
\[
u(x,t)=
\begin{cases}
\varphi(x-at), & x\ge at,
\\[0.4em]
g\!\left(t-\dfrac{x}{a}\right), & 0<x<at.
\end{cases}
\]

\textbf{Замечание о знаке $a$.} Если $a<0$, то характеристики выходят из области через $x=0$, а не входят в нее. В этом случае задавать произвольное граничное условие при $x=0$ нельзя: оно вообще не управляет решением внутри области так, как при $a>0$.

\paragraph{Мнемоническое правило / Суть на пальцах.}
На полупрямой всегда надо смотреть, откуда пришла характеристика: сверху из начального слоя $t=0$ или слева с границы $x=0$. Оттуда и берется значение решения.

\subsection*{1(e). Условие согласования в угловой точке}

\Answer
Для начально-краевой задачи на полупрямой или на отрезке начальные и граничные данные встречаются в \textbf{угловой точке} $(0,0)$. Чтобы решение было классическим, данные должны в этой точке не противоречить друг другу.

\textbf{Рассмотрим задачу}
\[
u_t+a u_x=f(x,t),
\qquad a>0,
\]
\[
u(x,0)=\varphi(x),
\qquad
u(0,t)=g(t).
\]

\textbf{Первое условие согласования.} Непрерывность решения в точке $(0,0)$ требует, чтобы значение, полученное из начального условия, совпадало со значением, полученным из граничного условия:
\[
\varphi(0)=g(0).
\]

\textbf{Почему это необходимо.}
Если бы
\[
\varphi(0)\ne g(0),
\]
то в одной и той же точке $(0,0)$ решение должно было бы принимать два разных значения. Тогда непрерывное, а тем более классическое, решение существовать не может.

\textbf{Второе условие согласования.} Если требуется, чтобы решение было не только непрерывным, но и непрерывно дифференцируемым, то нужно согласовать и первые производные. В пособии для уравнения
\[
u_t+a u_x=f(x,t)
\]
получается условие
\[
g'(0)+a\varphi'(0)=f(0,0).
\]

\textbf{Вывод этого условия.}

В точке $(0,0)$ должно выполняться само уравнение:
\[
u_t(0,0)+a u_x(0,0)=f(0,0).
\]

Но вдоль границы $x=0$ имеем
\[
u(0,t)=g(t),
\]
значит,
\[
u_t(0,0)=g'(0).
\]

А вдоль начальной прямой $t=0$ имеем
\[
u(x,0)=\varphi(x),
\]
значит,
\[
u_x(0,0)=\varphi'(0).
\]

Подставляя эти равенства в уравнение, получаем
\[
g'(0)+a\varphi'(0)=f(0,0).
\]

Для однородного уравнения $f\equiv 0$ это условие принимает вид
\[
g'(0)+a\varphi'(0)=0.
\]

\textbf{Смысл.} Условия согласования нужны для того, чтобы куски решения, построенные по начальным и граничным данным, можно было «склеить» вдоль особой характеристики, выходящей из угловой точки.

\paragraph{Мнемоническое правило / Суть на пальцах.}
В точке $(0,0)$ начальные и граничные данные должны не спорить друг с другом: сначала по значению, а если хотим гладкость, то и по первой производной. Иначе на углу получится излом или разрыв.

\subsection*{1(f). Полная производная вдоль характеристики и ее связь с уравнением переноса}

\Answer
Пусть характеристика задается уравнением
\[
x=x(t),
\qquad \frac{dx}{dt}=a.
\]
Тогда функция
\[
t\longmapsto u(x(t),t)
\]
является сложной функцией одной переменной.

\textbf{Определение.} \textbf{Полной производной вдоль характеристики} называется производная
\[
\frac{d}{dt}u(x(t),t).
\]

\textbf{Вычисление по правилу цепочки:}
\[
\frac{d}{dt}u(x(t),t)=u_t(x(t),t)+u_x(x(t),t)\frac{dx}{dt}.
\]

Так как
\[
\frac{dx}{dt}=a,
\]
получаем
\[
\frac{d}{dt}u(x(t),t)=u_t+a u_x.
\]

Следовательно, уравнение переноса
\[
u_t+a u_x=f(x,t)
\]
вдоль характеристики переписывается как обыкновенное дифференциальное уравнение
\[
\frac{d}{dt}u(x(t),t)=f(x(t),t).
\]

Это и есть главный смысл метода характеристик:
\begin{itemize}
  \item уравнение в частных производных заменяется на обыкновенное дифференциальное уравнение;
  \item для однородного уравнения $f\equiv 0$ получаем
  \[
  \frac{d}{dt}u(x(t),t)=0,
  \]
  то есть решение постоянно на характеристиках;
  \item для неоднородного уравнения
  \[
  \frac{d}{dt}u(x(t),t)=f(x(t),t),
  \]
  то есть вдоль характеристики решение изменяется ровно настолько, насколько его «подкачивает» правая часть $f$.
\end{itemize}

\textbf{Многомерное обобщение.} Для уравнения
\[
u_t+\mathbf{a}\cdot\nabla u=f
\]
при характеристиках
\[
\frac{d\mathbf{x}}{dt}=\mathbf{a}
\]
имеем
\[
\frac{d}{dt}u(\mathbf{x}(t),t)=u_t+\mathbf{a}\cdot\nabla u=f.
\]

\paragraph{Мнемоническое правило / Суть на пальцах.}
Полная производная вдоль характеристики отвечает на вопрос: «как меняется величина, если мы едем вместе с переносом?». Для однородного переноса ответ такой: никак не меняется.

\subsection*{1(g). Конвективная (субстанциональная, материальная) производная в гидродинамике}

\Answer
В гл.~5 пособия скорость жидкости обозначается вектором $\mathbf{u}=\mathbf{u}(x,t)$. Если $\psi=\psi(x,t)$ --- некоторая гидродинамическая величина (например, плотность, температура, соленость, концентрация), то \textbf{субстанциональной} или \textbf{материальной} производной называют скорость изменения $\psi$ вдоль траектории движущейся частицы жидкости.

\textbf{Определение.} Пусть траектория частицы задается системой
\[
\frac{d\mathbf{x}}{dt}=\mathbf{u}(\mathbf{x}(t),t).
\]
Тогда конвективная производная величины $\psi$ вдоль этой траектории равна
\[
\frac{D\psi}{Dt}=\frac{d}{dt}\psi(\mathbf{x}(t),t).
\]

\textbf{По правилу цепочки:}
\[
\frac{D\psi}{Dt}
=
\frac{\partial \psi}{\partial t}
+\mathbf{u}\cdot\nabla\psi.
\]

Тем самым субстанциональная производная распадается на две части:
\[
\frac{D\psi}{Dt}
=
\underbrace{\frac{\partial \psi}{\partial t}}_{\text{локальная часть}}
+\underbrace{\mathbf{u}\cdot\nabla\psi}_{\text{конвективная часть}}.
\]

\textbf{Смысл этих слагаемых:}
\begin{itemize}
  \item $\dfrac{\partial \psi}{\partial t}$ --- изменение в фиксированной точке пространства;
  \item $\mathbf{u}\cdot\nabla\psi$ --- изменение из-за того, что частица жидкости переходит в соседние точки, где значение $\psi$ иное.
\end{itemize}

\textbf{Связь с уравнением переноса.} Уравнение
\[
u_t+\mathbf{a}\cdot\nabla u=f
\]
есть точная многомерная модель того же механизма: вдоль траекторий, задаваемых полем скорости $\mathbf{a}$, выполняется
\[
\frac{Du}{Dt}=f.
\]
В гидродинамике в роли поля скорости выступает собственная скорость жидкости $\mathbf{u}$.

\textbf{Особенно важный частный случай.} Если в качестве $\psi$ взять сам вектор скорости жидкости $\mathbf{u}$, то
\[
\frac{D\mathbf{u}}{Dt}
=
\frac{\partial \mathbf{u}}{\partial t}
+(\mathbf{u}\cdot\nabla)\mathbf{u}.
\]
В пособии второе слагаемое
\[
(\mathbf{u}\cdot\nabla)\mathbf{u}
\]
называется \textbf{конвективным ускорением}. Тогда
\[
\frac{D\mathbf{u}}{Dt}
\]
есть полное ускорение частицы жидкости.

Именно в таком виде записывается уравнение Эйлера:
\[
\rho\frac{D\mathbf{u}}{Dt}=-\nabla p+\rho\mathbf{f}.
\]

\paragraph{Мнемоническое правило / Суть на пальцах.}
Локальная производная смотрит, как меняется поле в неподвижной точке. Субстанциональная --- как меняется поле для частицы, которую несет поток. Разница между ними и есть эффект переноса.

\section*{Практические задачи к разделу 1}

\subsection*{Задача 1(a). Уравнение \texorpdfstring{$u_t+2u_x=0$}{ut+2ux=0}}

\textbf{Дано.}
\[
u_t+2u_x=0.
\]

\textbf{Замечание о постановке.} В списке задач написано «решите задачу Коши», но начальное условие в формулировке не указано. Поэтому в строгом смысле задача Коши не замкнута. Ниже:
\begin{itemize}
  \item сначала находим \textbf{общее решение};
  \item затем записываем решение \textbf{корректно поставленной задачи Коши} при условии
  \[
  u(x,0)=\varphi(x).
  \]
\end{itemize}

\textbf{Теоретическое обоснование.}
Это линейное однородное уравнение переноса с постоянной скоростью
\[
a=2.
\]
Следовательно, применим метод характеристик.

\Solution
\textbf{Шаг 1. Находим характеристики.}

Они удовлетворяют уравнению
\[
\frac{dx}{dt}=2.
\]

Интегрируем:
\[
x=2t+C.
\]

Следовательно,
\[
x-2t=C.
\]

\textbf{Шаг 2. Вычисляем полную производную вдоль характеристики.}

Пусть $x=x(t)$ --- характеристика. Тогда
\[
\frac{d}{dt}u(x(t),t)=u_t+u_x\frac{dx}{dt}=u_t+2u_x.
\]

По уравнению
\[
u_t+2u_x=0,
\]
поэтому
\[
\frac{d}{dt}u(x(t),t)=0.
\]

\textbf{Комментарий.} Значит, вдоль каждой характеристики решение постоянно.

\textbf{Шаг 3. Записываем общее решение.}

Так как значение решения зависит только от номера характеристики $C=x-2t$, получаем
\[
u(x,t)=\Phi(x-2t),
\]
где $\Phi$ --- произвольная дифференцируемая функция.

\textbf{Шаг 4. Если задано начальное условие $u(x,0)=\varphi(x)$.}

Подставляем $t=0$:
\[
u(x,0)=\Phi(x)=\varphi(x).
\]

Значит,
\[
\Phi=\varphi.
\]

\textbf{Итак, решение задачи Коши имеет вид}
\[
u(x,t)=\varphi(x-2t).
\]

\textbf{Ответ.}
\begin{itemize}
  \item Общее решение:
  \[
  u(x,t)=\Phi(x-2t).
  \]
  \item Если понимать задачу как задачу Коши с условием $u(x,0)=\varphi(x)$, то
  \[
  u(x,t)=\varphi(x-2t).
  \]
\end{itemize}

\subsection*{Задача 1(b). Полная производная вдоль характеристики}

\textbf{Дано.}
\[
u_t+a u_x=0,
\qquad
u(x,0)=\varphi(x).
\]

\textbf{Теоретическое обоснование.}
Для уравнения переноса характеристики определяются уравнением
\[
\frac{dx}{dt}=a.
\]
Вдоль них уравнение должно перейти в обыкновенное дифференциальное уравнение.

\Solution
\textbf{Шаг 1. Находим характеристики.}

Интегрируем
\[
\frac{dx}{dt}=a:
\]
\[
x=at+C,
\qquad
x-at=C.
\]

\textbf{Шаг 2. Рассматриваем значение решения вдоль характеристики.}

Пусть
\[
x=x(t)=at+C.
\]
Тогда функция
\[
t\longmapsto u(x(t),t)
\]
зависит только от $t$.

\textbf{Шаг 3. Вычисляем полную производную.}

По правилу дифференцирования сложной функции
\[
\frac{d}{dt}u(x(t),t)=u_t(x(t),t)+u_x(x(t),t)\frac{dx}{dt}.
\]

Так как
\[
\frac{dx}{dt}=a,
\]
то
\[
\frac{d}{dt}u(x(t),t)=u_t+a u_x.
\]

\textbf{Шаг 4. Используем уравнение.}

Поскольку
\[
u_t+a u_x=0,
\]
то вдоль характеристики
\[
\frac{d}{dt}u(x(t),t)=0.
\]

\textbf{Комментарий.} Это означает, что значение $u$ не меняется при движении вдоль характеристики.

\textbf{Шаг 5. Записываем вывод о постоянстве решения.}

Если производная по $t$ равна нулю, то
\[
u(x(t),t)=\mathrm{const}
\]
на каждой характеристике.

Поскольку характеристика, проходящая через точку $(x,t)$, пересекает начальную прямую $t=0$ в точке
\[
(x-at,0),
\]
получаем
\[
u(x,t)=\varphi(x-at).
\]

\textbf{Ответ.}
\[
\frac{du}{dt}=u_t+a u_x.
\]
Для решения уравнения переноса
\[
u_t+a u_x=0
\]
вдоль характеристики выполняется
\[
\frac{du}{dt}=0,
\]
поэтому $u$ постоянно на характеристиках.

\section*{2. Метод энергетических неравенств}

\textbf{Обозначение норм.} В этом разделе для согласования с практической задачей будем рассматривать, если не оговорено иное, функцию
\[
u=u(x,t), \qquad x>0, \quad 0<t\le T,
\]
и пользоваться обозначениями
\[
(u,v)=\int_0^{+\infty}u(x,t)v(x,t)\,dx,
\qquad
\|u(\cdot,t)\|^2=(u,u)=\int_0^{+\infty}u^2(x,t)\,dx.
\]
Предполагается, что решение достаточно гладко, а функции $u(x,t)$ и $u_x(x,t)$ убывают при $x\to+\infty$ настолько быстро, чтобы все интегрирования по частям были корректны.

\subsection*{2(a). Суть энергетического метода}

\Answer
\textbf{Энергетический метод} --- это метод получения \textbf{априорных оценок} решения с помощью самого дифференциального уравнения, без использования явной формулы решения.

В логике пособия Алексеев подчёркивает две ключевые мысли.
\begin{itemize}
  \item Мы не пытаемся сначала найти точное решение, а выводим неравенства, которым \textbf{обязано} удовлетворять любое достаточно гладкое решение.
  \item Полученные оценки затем используются для доказательства \textbf{единственности}, \textbf{устойчивости}, а часто и \textbf{существования} решения.
\end{itemize}

\textbf{Главная идея метода.} Нужно выбрать такую функциональную величину, которая:
\begin{itemize}
  \item естественно связана с уравнением;
  \item неотрицательна;
  \item позволяет после умножения уравнения на подходящую функцию и интегрирования получить полезное тождество.
\end{itemize}

Для уравнений теплопроводности, переноса и колебаний такой величиной обычно служит квадрат нормы решения:
\[
E(t)=\|u(\cdot,t)\|^2.
\]
Эта величина и называется \textbf{энергией} системы в математическом смысле.

\textbf{Почему метод называется энергетическим.}
\begin{itemize}
  \item В очень многих физических моделях величина
  \[
  \|u(\cdot,t)\|^2
  \]
  либо сама является энергией, либо пропорциональна энергии.
  \item Даже когда прямой физический смысл не буквальный, эта величина играет ту же роль: она измеряет «размер» состояния системы и позволяет контролировать его эволюцию во времени.
\end{itemize}

\textbf{Типичная схема энергетического метода:}
\begin{enumerate}
  \item Берут уравнение и умножают его на $u$ или на другую специально выбранную функцию.
  \item Интегрируют по пространственной переменной или по области.
  \item Переносят производные с помощью интегрирования по частям.
  \item Получают \textbf{энергетическое тождество} или \textbf{энергетическое неравенство}.
  \item Из него выводят \textbf{априорную оценку}.
  \item Применяют эту оценку к разности двух решений, чтобы получить \textbf{единственность} и \textbf{устойчивость}.
\end{enumerate}

\textbf{Что значит априорная оценка.} Это оценка, полученная \textbf{до} нахождения точного решения и только из структуры уравнения и исходных данных.

\paragraph{Мнемоническое правило / Суть на пальцах.}
Энергетический метод не спрашивает: «каково решение в явном виде?». Он спрашивает: «насколько большим решение вообще может стать, если уравнение и данные уже зафиксированы?». Мы оцениваем не саму траекторию, а запас её «энергии».

\subsection*{2(b). Энергетическое тождество для \texorpdfstring{$u_t-a u_{xx}+\gamma u=f$}{ut-a uxx+gamma u=f}}

\Answer
Рассмотрим уравнение
\[
u_t-a u_{xx}+\gamma u=f(x,t),
\qquad a>0,\quad \gamma\ge 0,
\]
на полупрямой $x>0$ при условиях
\[
u(0,t)=0,
\qquad
u(x,0)=\varphi(x),
\]
и предположим, что
\[
u(x,t)\to 0,
\qquad
u_x(x,t)\to 0,
\qquad x\to +\infty.
\]

\textbf{Цель.} Нужно получить тождество, связывающее норму решения, норму производной $u_x$ и правую часть $f$.

\textbf{Шаг 1.} Умножаем уравнение на $2u$:
\[
2u u_t-2a u u_{xx}+2\gamma u^2=2fu.
\]

\textbf{Комментарий.} Множитель $2u$ выбран потому, что тогда первое слагаемое превратится в производную от $u^2$, а значит --- в производную от квадрата нормы.

\textbf{Шаг 2.} Интегрируем по $x$ на $(0,+\infty)$:
\[
2\int_0^{+\infty}u u_t\,dx
-2a\int_0^{+\infty}u u_{xx}\,dx
+2\gamma\int_0^{+\infty}u^2\,dx
=
2\int_0^{+\infty}fu\,dx.
\]

\textbf{Шаг 3.} Преобразуем первое слагаемое:
\[
2\int_0^{+\infty}u u_t\,dx
=
\frac{d}{dt}\int_0^{+\infty}u^2\,dx
=
\frac{d}{dt}\|u(\cdot,t)\|^2.
\]

\textbf{Шаг 4.} Преобразуем второе слагаемое интегрированием по частям:
\[
\int_0^{+\infty}u u_{xx}\,dx
=
\left.u u_x\right|_{0}^{+\infty}
-\int_0^{+\infty}u_x^2\,dx.
\]

По условиям
\[
u(0,t)=0,
\qquad
u(x,t)\to 0,
\qquad
u_x(x,t)\to 0
\quad (x\to +\infty),
\]
граничный член исчезает:
\[
\left.u u_x\right|_{0}^{+\infty}=0.
\]

Следовательно,
\[
\int_0^{+\infty}u u_{xx}\,dx=-\int_0^{+\infty}u_x^2\,dx,
\]
а потому
\[
-2a\int_0^{+\infty}u u_{xx}\,dx
=
2a\int_0^{+\infty}u_x^2\,dx.
\]

\textbf{Шаг 5.} Собираем всё вместе:
\[
\frac{d}{dt}\|u(\cdot,t)\|^2
+2a\|u_x(\cdot,t)\|^2
+2\gamma\|u(\cdot,t)\|^2
=
2(f(\cdot,t),u(\cdot,t)).
\]

Это и есть \textbf{дифференциальное энергетическое тождество}.

\textbf{Шаг 6.} Интегрируем по времени на отрезке $[0,t]$:
\[
\int_0^t \frac{d}{d\tau}\|u(\cdot,\tau)\|^2\,d\tau
+2a\int_0^t \|u_x(\cdot,\tau)\|^2\,d\tau
+2\gamma\int_0^t \|u(\cdot,\tau)\|^2\,d\tau
=
2\int_0^t (f(\cdot,\tau),u(\cdot,\tau))\,d\tau.
\]

Первый интеграл вычисляется по формуле Ньютона--Лейбница:
\[
\int_0^t \frac{d}{d\tau}\|u(\cdot,\tau)\|^2\,d\tau
=
\|u(\cdot,t)\|^2-\|u(\cdot,0)\|^2.
\]

Так как
\[
u(x,0)=\varphi(x),
\]
получаем
\[
\|u(\cdot,0)\|^2=\|\varphi\|^2.
\]

Следовательно, \textbf{интегральное энергетическое тождество} имеет вид
\[
\|u(\cdot,t)\|^2
+2a\int_0^t \|u_x(\cdot,\tau)\|^2\,d\tau
+2\gamma\int_0^t \|u(\cdot,\tau)\|^2\,d\tau
=
\|\varphi\|^2
+2\int_0^t (f(\cdot,\tau),u(\cdot,\tau))\,d\tau.
\]

\paragraph{Мнемоническое правило / Суть на пальцах.}
Уравнение теплопроводности слагает «баланс энергии»: текущая энергия равна начальной энергии плюс работа источника $f$ минус то, что рассеивается из-за сглаживания, то есть через $u_x$.

\subsection*{2(c). Как из энергетического тождества получается априорная оценка}

\Answer
Исходной точкой служит дифференциальное энергетическое тождество
\[
\frac{d}{dt}\|u(\cdot,t)\|^2
+2a\|u_x(\cdot,t)\|^2
+2\gamma\|u(\cdot,t)\|^2
=
2(f(\cdot,t),u(\cdot,t)).
\]

\textbf{Шаг 1.} Отбрасываем неотрицательные слагаемые в левой части:
\[
2a\|u_x(\cdot,t)\|^2\ge 0,
\qquad
2\gamma\|u(\cdot,t)\|^2\ge 0,
\]
так как
\[
a>0,
\qquad
\gamma\ge 0.
\]

Получаем неравенство
\[
\frac{d}{dt}\|u(\cdot,t)\|^2
\le
2(f(\cdot,t),u(\cdot,t)).
\]

\textbf{Шаг 2.} Применяем неравенство Коши--Буняковского:
\[
|(f(\cdot,t),u(\cdot,t))|
\le
\|f(\cdot,t)\|\,\|u(\cdot,t)\|.
\]

Значит,
\[
\frac{d}{dt}\|u(\cdot,t)\|^2
\le
2\|f(\cdot,t)\|\,\|u(\cdot,t)\|.
\]

\textbf{Шаг 3.} Обозначим
\[
y(t)=\|u(\cdot,t)\|.
\]
Тогда
\[
\frac{d}{dt}y^2(t)=2y(t)y'(t)
\]
во всех точках, где $y(t)>0$. Следовательно,
\[
2y(t)y'(t)\le 2\|f(\cdot,t)\|\,y(t).
\]

Если $y(t)>0$, можно сократить на $2y(t)$:
\[
y'(t)\le \|f(\cdot,t)\|.
\]

Если же $y(t)=0$, то итоговая оценка тем более сохраняется по непрерывности. Поэтому можно писать в интегральном смысле
\[
y'(t)\le \|f(\cdot,t)\|.
\]

\textbf{Шаг 4.} Интегрируем по времени от $0$ до $t$:
\[
y(t)-y(0)\le \int_0^t \|f(\cdot,\tau)\|\,d\tau.
\]

\textbf{Шаг 5.} Подставляем
\[
y(0)=\|u(\cdot,0)\|=\|\varphi\|.
\]
Получаем априорную оценку
\[
\|u(\cdot,t)\|
\le
\|\varphi\|
+\int_0^t \|f(\cdot,\tau)\|\,d\tau.
\]

\textbf{Именно эта оценка и требуется в практической задаче.}

\paragraph{Мнемоническое правило / Суть на пальцах.}
Чтобы оценить решение сверху, мы просто забываем про «полезное» затухание в левой части и оставляем только то, чем правую часть можно раскачать систему. Получается грубая, но очень надёжная оценка.

\subsection*{2(d). Единственность и устойчивость из априорной оценки}

\Answer
Пусть для уравнения
\[
u_t-a u_{xx}+\gamma u=f
\]
с заданными начальными и граничными условиями получена априорная оценка
\[
\|u(\cdot,t)\|
\le
\|\varphi\|
+\int_0^t \|f(\cdot,\tau)\|\,d\tau.
\]

\textbf{1. Единственность.}

Пусть $u_1$ и $u_2$ --- два решения одной и той же задачи с одинаковыми данными. Рассмотрим их разность
\[
w=u_1-u_2.
\]

Тогда, вычитая одно уравнение из другого, получаем
\[
w_t-a w_{xx}+\gamma w=0.
\]

Для начального условия имеем
\[
w(x,0)=u_1(x,0)-u_2(x,0)=\varphi(x)-\varphi(x)=0.
\]

Для граничного условия:
\[
w(0,t)=u_1(0,t)-u_2(0,t)=0.
\]

Применяем априорную оценку к функции $w$:
\[
\|w(\cdot,t)\|
\le
\|w(\cdot,0)\|
+\int_0^t \|0\|\,d\tau=0.
\]

Следовательно,
\[
\|w(\cdot,t)\|=0,
\]
то есть
\[
w(x,t)\equiv 0.
\]

Значит,
\[
u_1(x,t)\equiv u_2(x,t).
\]

Итак, решение единственно.

\textbf{2. Устойчивость.}

Пусть теперь $u_1$ и $u_2$ отвечают разным данным
\[
(\varphi_1,f_1),
\qquad
(\varphi_2,f_2).
\]

Снова положим
\[
w=u_1-u_2.
\]
Тогда
\[
w_t-a w_{xx}+\gamma w=f_1-f_2,
\]
\[
w(x,0)=\varphi_1(x)-\varphi_2(x),
\qquad
w(0,t)=0
\]
в случае одинакового граничного условия.

Применяя априорную оценку к $w$, получаем
\[
\|u_1(\cdot,t)-u_2(\cdot,t)\|
\le
\|\varphi_1-\varphi_2\|
+\int_0^t \|f_1(\cdot,\tau)-f_2(\cdot,\tau)\|\,d\tau.
\]

\textbf{Смысл этой оценки.} Малые изменения исходных данных вызывают малые изменения решения. Это и есть устойчивость.

\paragraph{Мнемоническое правило / Суть на пальцах.}
Для единственности надо вычесть два решения и показать, что разность обязана быть нулём. Для устойчивости делаем то же самое, только теперь разность данных не нулевая, а малая, и потому мала и разность решений.

\subsection*{2(e). Многомерный случай}

\Answer
В многомерном случае вместо уравнения
\[
u_t-a u_{xx}+\gamma u=f
\]
рассматривают уравнение
\[
u_t-a\Delta u+\gamma u=f(x,t),
\qquad x\in\Omega\subset\R^n,
\]
с, например, однородным условием Дирихле
\[
u|_{\partial\Omega}=0.
\]

\textbf{Смысл энергетического метода не меняется.} Мы снова умножаем уравнение на $2u$ и интегрируем по области $\Omega$:
\[
2\int_\Omega u u_t\,dx
-2a\int_\Omega u\Delta u\,dx
+2\gamma\int_\Omega u^2\,dx
=
2\int_\Omega fu\,dx.
\]

\textbf{Первое слагаемое:}
\[
2\int_\Omega u u_t\,dx=\frac{d}{dt}\int_\Omega u^2\,dx.
\]

\textbf{Второе слагаемое:} применяем формулу Грина:
\[
\int_\Omega u\Delta u\,dx
=
\int_{\partial\Omega}u\frac{\partial u}{\partial n}\,dS
-\int_\Omega |\nabla u|^2\,dx.
\]

Так как при условии Дирихле
\[
u|_{\partial\Omega}=0,
\]
граничный интеграл исчезает, и потому
\[
-2a\int_\Omega u\Delta u\,dx
=
2a\int_\Omega |\nabla u|^2\,dx.
\]

Следовательно, получаем многомерное энергетическое тождество
\[
\frac{d}{dt}\|u(\cdot,t)\|_{L^2(\Omega)}^2
+2a\|\nabla u(\cdot,t)\|_{L^2(\Omega)}^2
+2\gamma\|u(\cdot,t)\|_{L^2(\Omega)}^2
=
2(f(\cdot,t),u(\cdot,t))_{L^2(\Omega)}.
\]

После интегрирования по времени:
\[
\begin{aligned}
\|u(\cdot,t)\|_{L^2(\Omega)}^2
&+2a\int_0^t \|\nabla u(\cdot,\tau)\|_{L^2(\Omega)}^2\,d\tau
+2\gamma\int_0^t \|u(\cdot,\tau)\|_{L^2(\Omega)}^2\,d\tau
\\
&=
\|u(\cdot,0)\|_{L^2(\Omega)}^2
+2\int_0^t (f(\cdot,\tau),u(\cdot,\tau))_{L^2(\Omega)}\,d\tau.
\end{aligned}
\]

Далее всё идёт так же, как в одномерном случае:
\begin{itemize}
  \item отбрасываем неотрицательные слагаемые;
  \item применяем неравенство Коши--Буняковского;
  \item получаем априорную оценку;
  \item затем доказываем единственность и устойчивость.
\end{itemize}

\paragraph{Мнемоническое правило / Суть на пальцах.}
В многомерном случае ничего принципиально нового не происходит: $u_x^2$ просто заменяется на $|\nabla u|^2$. То есть «шероховатость» решения измеряется уже не одной производной, а всем градиентом.

\section*{Практические задачи к разделу 2}

\subsection*{Задача 2(a). Вывести априорную оценку}

\textbf{Дано.}

Рассматривается уравнение
\[
u_t-a u_{xx}+\gamma u=f,
\qquad \gamma\ge 0,
\]
с условиями
\[
u(0,t)=0,
\qquad
u(x,0)=\varphi(x).
\]

\textbf{Что нужно получить:}
\[
\|u(\cdot,t)\|
\le
\|\varphi\|
+\int_0^t \|f(\cdot,\tau)\|\,d\tau.
\]

\textbf{Теоретическое обоснование.}
Для линейного параболического уравнения второго порядка естественно применять энергетический метод: умножаем уравнение на решение $u$, интегрируем по пространственной переменной и затем оцениваем правую часть через неравенство Коши--Буняковского.

\Solution
\textbf{Шаг 1. Умножаем уравнение на $2u$.}

Получаем
\[
2u u_t-2a u u_{xx}+2\gamma u^2=2fu.
\]

\textbf{Комментарий.} Так мы готовим выражение к интегрированию: первое слагаемое станет производной от нормы.

\textbf{Шаг 2. Интегрируем по $x$ на полупрямой $(0,+\infty)$.}

\[
2\int_0^{+\infty}u u_t\,dx
-2a\int_0^{+\infty}u u_{xx}\,dx
+2\gamma\int_0^{+\infty}u^2\,dx
=
2\int_0^{+\infty}fu\,dx.
\]

\textbf{Шаг 3. Преобразуем первое слагаемое.}

\[
2\int_0^{+\infty}u u_t\,dx
=
\frac{d}{dt}\int_0^{+\infty}u^2\,dx
=
\frac{d}{dt}\|u(\cdot,t)\|^2.
\]

\textbf{Шаг 4. Интегрируем по частям слагаемое с $u_{xx}$.}

\[
\int_0^{+\infty}u u_{xx}\,dx
=
\left.u u_x\right|_0^{+\infty}
-\int_0^{+\infty}u_x^2\,dx.
\]

Из граничного условия $u(0,t)=0$ и условий убывания на бесконечности следует
\[
\left.u u_x\right|_0^{+\infty}=0.
\]

Поэтому
\[
-2a\int_0^{+\infty}u u_{xx}\,dx
=
2a\int_0^{+\infty}u_x^2\,dx\ge 0.
\]

\textbf{Шаг 5. Получаем энергетическое тождество.}

\[
\frac{d}{dt}\|u(\cdot,t)\|^2
+2a\|u_x(\cdot,t)\|^2
+2\gamma\|u(\cdot,t)\|^2
=
2(f(\cdot,t),u(\cdot,t)).
\]

\textbf{Шаг 6. Отбрасываем неотрицательные слагаемые слева.}

Так как
\[
2a\|u_x(\cdot,t)\|^2\ge 0,
\qquad
2\gamma\|u(\cdot,t)\|^2\ge 0,
\]
то
\[
\frac{d}{dt}\|u(\cdot,t)\|^2
\le
2(f(\cdot,t),u(\cdot,t)).
\]

\textbf{Шаг 7. Применяем неравенство Коши--Буняковского.}

\[
|(f(\cdot,t),u(\cdot,t))|
\le
\|f(\cdot,t)\|\,\|u(\cdot,t)\|.
\]

Значит,
\[
\frac{d}{dt}\|u(\cdot,t)\|^2
\le
2\|f(\cdot,t)\|\,\|u(\cdot,t)\|.
\]

\textbf{Шаг 8. Вводим обозначение $y(t)=\|u(\cdot,t)\|$.}

Тогда
\[
\frac{d}{dt}y^2(t)\le 2\|f(\cdot,t)\|\,y(t).
\]

При $y(t)>0$ делим на $2y(t)$:
\[
y'(t)\le \|f(\cdot,t)\|.
\]

\textbf{Шаг 9. Интегрируем по времени от $0$ до $t$.}

\[
y(t)-y(0)\le \int_0^t \|f(\cdot,\tau)\|\,d\tau.
\]

Но
\[
y(0)=\|u(\cdot,0)\|=\|\varphi\|.
\]

Следовательно,
\[
\|u(\cdot,t)\|
\le
\|\varphi\|
+\int_0^t \|f(\cdot,\tau)\|\,d\tau.
\]

\textbf{Ответ.}
\[
\|u(\cdot,t)\|
\le
\|\varphi\|
+\int_0^t \|f(\cdot,\tau)\|\,d\tau.
\]

\subsection*{Задача 2(b). Доказать единственность с помощью энергетического тождества}

\textbf{Дано.}

Та же задача:
\[
u_t-a u_{xx}+\gamma u=f,
\qquad \gamma\ge 0,
\]
\[
u(0,t)=0,
\qquad
u(x,0)=\varphi(x).
\]

\textbf{Нужно доказать.} Решение единственно.

\textbf{Теоретическое обоснование.}
Для линейных задач единственность обычно доказывают переходом к разности двух решений. Тогда возникает однородная задача, к которой применяется энергетическое тождество.

\Solution
\textbf{Шаг 1. Предположим, что существует два решения.}

Пусть $u_1$ и $u_2$ --- два решения одной и той же задачи.

\textbf{Шаг 2. Рассмотрим их разность.}

Положим
\[
w=u_1-u_2.
\]

\textbf{Шаг 3. Выведем уравнение для $w$.}

Так как оба решения удовлетворяют одному и тому же уравнению
\[
u_t-a u_{xx}+\gamma u=f,
\]
то после вычитания получаем
\[
w_t-a w_{xx}+\gamma w=0.
\]

\textbf{Шаг 4. Выпишем условия для $w$.}

Из одинаковости начальных условий следует
\[
w(x,0)=u_1(x,0)-u_2(x,0)=\varphi(x)-\varphi(x)=0.
\]

Из одинаковости граничных условий следует
\[
w(0,t)=u_1(0,t)-u_2(0,t)=0.
\]

\textbf{Шаг 5. Применяем энергетическое тождество к функции $w$.}

Поскольку правая часть равна нулю, имеем
\[
\frac{d}{dt}\|w(\cdot,t)\|^2
+2a\|w_x(\cdot,t)\|^2
+2\gamma\|w(\cdot,t)\|^2
=0.
\]

\textbf{Шаг 6. Интегрируем по времени от $0$ до $t$.}

\[
\|w(\cdot,t)\|^2
+2a\int_0^t \|w_x(\cdot,\tau)\|^2\,d\tau
+2\gamma\int_0^t \|w(\cdot,\tau)\|^2\,d\tau
=
\|w(\cdot,0)\|^2.
\]

Но
\[
\|w(\cdot,0)\|^2=0,
\]
так как $w(x,0)=0$.

\textbf{Шаг 7. Используем неотрицательность всех слагаемых.}

Все три слагаемых в левой части неотрицательны. Их сумма равна нулю, следовательно, каждое из них равно нулю. В частности,
\[
\|w(\cdot,t)\|^2=0.
\]

Значит,
\[
w(x,t)\equiv 0.
\]

Следовательно,
\[
u_1(x,t)\equiv u_2(x,t).
\]

\textbf{Ответ.}
Разность двух решений удовлетворяет однородной задаче и имеет нулевые начальные и граничные данные. Энергетическое тождество даёт
\[
\|w(\cdot,t)\|^2=0,
\]
откуда
\[
w\equiv 0.
\]
Значит, решение единственно.

\section*{3. Уравнение Бюргерса и градиентная катастрофа}

\textit{Замечание по источнику. В используемом пособии Г.\,В.~Алексеева уравнение Бюргерса не выделено как отдельный параграф. Поэтому ниже оно рассматривается как специальный случай квазилинейного уравнения первого порядка из \S 2.6 гл.~2, а терминология «градиентная катастрофа», «ударная волна» и «вязкий/невязкий Бюргерс» берётся из приложенного списка вопросов.}

\subsection*{3(a). Вязкое и невязкое уравнение Бюргерса. Отличие от линейного переноса}

\Answer
\textbf{Невязкое уравнение Бюргерса} имеет вид
\[
u_t+u\,u_x=0.
\]

\textbf{Вязкое уравнение Бюргерса} имеет вид
\[
u_t+u\,u_x=\nu u_{xx},
\qquad \nu>0.
\]

\textbf{Связь с квазилинейным уравнением первого порядка.}
В гл.~2 пособия рассматривается квазилинейное уравнение
\[
a(x,y,u)\,u_x+b(x,y,u)\,u_y=f(x,y,u).
\]
Если положить
\[
y=t,\qquad a=u,\qquad b=1,\qquad f=0,
\]
то получим как раз невязкое уравнение Бюргерса
\[
u_t+u\,u_x=0.
\]
Следовательно, оно является \textbf{квазилинейным уравнением первого порядка}.

\textbf{Отличие от линейного уравнения переноса.}

Линейное уравнение переноса имеет вид
\[
u_t+a u_x=0,
\qquad a=\mathrm{const}
\]
или, более общо,
\[
u_t+a(x,t)u_x=0,
\]
где скорость переноса \textbf{не зависит от самого решения}.

В уравнении Бюргерса ситуация иная:
\[
u_t+u\,u_x=0.
\]
Здесь коэффициент при $u_x$ равен самому решению $u$. Поэтому:
\begin{itemize}
  \item скорость переноса зависит от величины, которая переносится;
  \item характеристики определяются не заранее, а вместе с решением;
  \item разные участки профиля движутся с разными скоростями;
  \item возможны пересечения характеристик и потеря гладкости.
\end{itemize}

\textbf{Роль вязкого члена.}
В вязком уравнении Бюргерса появляется слагаемое
\[
\nu u_{xx},
\]
которое отвечает за диффузионное сглаживание решения. Поэтому вязкое уравнение Бюргерса уже не является уравнением первого порядка: это \textbf{нелинейное параболическое уравнение второго порядка}.

\paragraph{Мнемоническое правило / Суть на пальцах.}
Линейный перенос --- это «весь профиль едет с одной и той же скоростью». Невязкий Бюргерс --- это «каждая точка профиля едет со своей собственной скоростью, равной её высоте». Вязкий Бюргерс --- то же самое, но с постоянным подглаживанием.

\subsection*{3(b). Градиентная катастрофа (опрокидывание волны)}

\Answer
\textbf{Градиентной катастрофой} называют ситуацию, когда решение $u(x,t)$ остаётся ограниченным, но его производная по пространству
\[
u_x(x,t)
\]
обращается в бесконечность за конечное время. Геометрически это означает, что профиль решения становится вертикальным, после чего классическое гладкое решение перестаёт существовать.

\textbf{Рассмотрим задачу Коши}
\[
u_t+u\,u_x=0,
\qquad
u(x,0)=\varphi(x).
\]

\textbf{Шаг 1.} Для невязкого уравнения Бюргерса характеристики определяются системой
\[
\frac{dx}{dt}=u,
\qquad
\frac{du}{dt}=0.
\]

\textbf{Комментарий.} Второе уравнение означает, что вдоль характеристики значение решения сохраняется.

\textbf{Шаг 2.} Если характеристика выходит из точки $x=\xi$ начальной прямой $t=0$, то
\[
u=\varphi(\xi)
\]
на всей этой характеристике. Поэтому первое характеристическое уравнение принимает вид
\[
\frac{dx}{dt}=\varphi(\xi).
\]

\textbf{Шаг 3.} Интегрируем по $t$:
\[
x=\xi+t\varphi(\xi).
\]

Таким образом, классическое решение задаётся параметрически формулами
\[
u(x,t)=\varphi(\xi),
\qquad
x=\xi+t\varphi(\xi).
\]

\textbf{Шаг 4.} Чтобы решение было однозначным гладким, отображение
\[
\xi\longmapsto x(\xi,t)=\xi+t\varphi(\xi)
\]
должно быть взаимно однозначным. Для этого необходимо, чтобы
\[
\frac{\partial x}{\partial \xi}=1+t\varphi'(\xi)\ne 0.
\]

\textbf{Шаг 5.} Пока
\[
1+t\varphi'(\xi)>0
\]
для всех $\xi$, характеристики не пересекаются, и классическое решение существует.

\textbf{Шаг 6.} Если же найдётся точка, где
\[
1+t\varphi'(\xi)=0,
\]
то две соседние характеристики слипаются. В этот момент и возникает градиентная катастрофа.

\textbf{Формула для производной $u_x$.}
Дифференцируя соотношения
\[
u=\varphi(\xi),
\qquad
x=\xi+t\varphi(\xi),
\]
по параметру $\xi$, получаем
\[
\frac{\partial u}{\partial \xi}=\varphi'(\xi),
\qquad
\frac{\partial x}{\partial \xi}=1+t\varphi'(\xi).
\]
Следовательно,
\[
u_x=\frac{u_\xi}{x_\xi}
=
\frac{\varphi'(\xi)}{1+t\varphi'(\xi)}.
\]

Если
\[
1+t\varphi'(\xi)\to 0,
\]
то
\[
|u_x|\to\infty.
\]
При этом сама величина
\[
u=\varphi(\xi)
\]
остаётся ограниченной, если ограничена начальная функция $\varphi$.

\textbf{Время разрушения гладкого решения.}
Если
\[
\min_{\xi\in\R}\varphi'(\xi)<0,
\]
то первое время возникновения градиентной катастрофы равно
\[
t_*=
-\frac{1}{\min\limits_{\xi}\varphi'(\xi)}.
\]

Если же
\[
\varphi'(\xi)\ge 0
\quad \text{для всех } \xi,
\]
то характеристики не пересекаются, и механизм опрокидывания не возникает.

\textbf{Физическая причина.}
Участки профиля, где $u$ больше, движутся быстрее. Если более высокие участки находятся сзади, они начинают догонять более низкие передние участки. Из-за этого фронт «заламывается», и возникает опрокидывание волны.

\paragraph{Мнемоническое правило / Суть на пальцах.}
Градиентная катастрофа --- это момент, когда «быстрые» куски волны догнали «медленные». Значения ещё конечны, но наклон профиля уже уходит в бесконечность.

\subsection*{3(c). Характеристики для уравнения Бюргерса и их пересечение}

\Answer
Для невязкого уравнения Бюргерса
\[
u_t+u\,u_x=0
\]
характеристики ищутся так же, как для общего квазилинейного уравнения первого порядка.

\textbf{Шаг 1.} Запишем систему характеристик:
\[
\frac{dx}{dt}=u,
\qquad
\frac{du}{dt}=0.
\]

\textbf{Шаг 2.} Из второго уравнения следует, что вдоль каждой характеристики
\[
u=\mathrm{const}.
\]
Если характеристика стартует из точки $x=\xi$ начальной прямой $t=0$, то
\[
u=\varphi(\xi).
\]

\textbf{Шаг 3.} Подставляем это в первое уравнение:
\[
\frac{dx}{dt}=\varphi(\xi).
\]
Интегрируя, получаем
\[
x=\xi+t\varphi(\xi).
\]

Следовательно, \textbf{характеристика, выходящая из точки $\xi$}, задаётся формулой
\[
x=\xi+t\varphi(\xi).
\]

\textbf{Почему характеристики могут пересекаться.}
Для линейного переноса
\[
u_t+a u_x=0
\]
все характеристики имеют один и тот же наклон, так как
\[
\frac{dx}{dt}=a.
\]
Поэтому они параллельны и не пересекаются.

У уравнения Бюргерса наклон характеристики зависит от начального значения:
\[
\frac{dx}{dt}=\varphi(\xi).
\]
Если для разных $\xi$ значения $\varphi(\xi)$ различны, то и скорости различных характеристик различны.

\textbf{Условие пересечения.}
Характеристики пересекаются тогда, когда отображение
\[
\xi\mapsto x(\xi,t)=\xi+t\varphi(\xi)
\]
теряет монотонность, то есть когда
\[
x_\xi(\xi,t)=1+t\varphi'(\xi)=0.
\]

\textbf{Геометрический смысл.}
\begin{itemize}
  \item если $x_\xi>0$, то разные характеристики идут отдельно;
  \item если $x_\xi=0$, две соседние характеристики касаются;
  \item если $x_\xi<0$, параметрическое задание перестаёт быть однозначным, и классическое решение уже невозможно.
\end{itemize}

\textbf{Вывод.}
Пересечение характеристик для уравнения Бюргерса --- следствие нелинейности: скорость переноса определяется самим решением.

\paragraph{Мнемоническое правило / Суть на пальцах.}
У Бюргерса характеристики не параллельны: чем выше волна, тем круче и быстрее идёт её характеристика. Из-за этого они могут «сойтись в одну точку».

\subsection*{3(d). Как вязкость предотвращает градиентную катастрофу}

\Answer
Вязкое уравнение Бюргерса имеет вид
\[
u_t+u\,u_x=\nu u_{xx},
\qquad \nu>0.
\]

\textbf{Главное отличие от невязкого случая.}
В невязком уравнении
\[
u_t+u\,u_x=0
\]
механизм эволюции чисто транспортный: каждая точка профиля просто переносится со скоростью, равной её высоте. Из-за этого и возможно пересечение характеристик.

В вязком случае появляется диффузионное слагаемое
\[
\nu u_{xx},
\]
которое:
\begin{itemize}
  \item сглаживает резкие перепады;
  \item переносит информацию между соседними точками;
  \item препятствует бесконечному росту градиента.
\end{itemize}

\textbf{Качественный смысл члена $\nu u_{xx}$.}
Если в некоторой точке профиль слишком выпуклый или слишком вогнутый, вторая производная велика по модулю, и диффузионный член начинает активно выравнивать профиль. Поэтому фронт не успевает стать вертикальным: прежде чем возникнет бесконечный наклон, вязкость его размывает.

\textbf{С математической точки зрения.}
Уравнение
\[
u_t+u\,u_x=\nu u_{xx}
\]
является нелинейным параболическим уравнением. Для параболических уравнений характерен сглаживающий эффект: гладкие начальные данные порождают гладкие решения, и для $t>0$ решение становится более регулярным, а не менее.

\textbf{Физическая роль вязкости.}
Вязкость моделирует внутреннее трение среды. Она:
\begin{itemize}
  \item рассеивает энергию мелкомасштабных градиентов;
  \item сглаживает крутые фронты;
  \item переводит резкий разрыв в узкий, но гладкий переходный слой.
\end{itemize}

\textbf{Итог.}
Вязкий член не отменяет стремление нелинейного переноса к «схлопыванию» фронта, но постоянно конкурирует с ним. В результате вместо мгновенного опрокидывания возникает гладкий фронт конечной толщины.

\paragraph{Мнемоническое правило / Суть на пальцах.}
Нелинейный перенос хочет сделать фронт круче. Вязкость всё время его подглаживает. Если невязкий Бюргерс «ломает» волну, то вязкий Бюргерс её «закругляет».

\subsection*{3(e). Ударная волна и её связь с градиентной катастрофой}

\Answer
Невязкое уравнение Бюргерса можно записать в \textbf{законе сохранения}
\[
u_t+\left(\frac{u^2}{2}\right)_x=0.
\]

\textbf{Ударная волна} --- это разрывное решение такого уравнения, у которого функция $u$ испытывает скачок на некоторой движущейся кривой
\[
x=s(t).
\]

\textbf{Почему возникает разрыв.}
Пока характеристики не пересекаются, существует классическое гладкое решение. Но после момента градиентной катастрофы классическое решение разрушается: в одной и той же точке плоскости $(x,t)$ характеристический метод начинает давать несколько значений. Чтобы сохранить физически осмысленное решение, переходят к \textbf{обобщённым} или \textbf{слабым} решениям. В их классе и появляются разрывы --- ударные волны.

\textbf{Скорость движения разрыва.}
Для закона сохранения
\[
u_t+\bigl(F(u)\bigr)_x=0,
\qquad
F(u)=\frac{u^2}{2},
\]
условие Ранкина--Гюгонио даёт
\[
s'(t)=\frac{F(u_-)-F(u_+)}{u_--u_+},
\]
где $u_-$ и $u_+$ --- левые и правые пределы решения на разрыве.

Подставляя
\[
F(u)=\frac{u^2}{2},
\]
получаем
\[
s'(t)=\frac{\dfrac{u_-^2}{2}-\dfrac{u_+^2}{2}}{u_--u_+}
=
\frac{u_-+u_+}{2}.
\]

\textbf{Связь с градиентной катастрофой.}
\begin{itemize}
  \item сначала гладкое решение становится всё круче;
  \item затем в конечный момент времени возникает градиентная катастрофа;
  \item после этого гладкое решение заменяется слабым решением с разрывом;
  \item этот разрыв и интерпретируется как ударная волна.
\end{itemize}

\textbf{Связь с вязким уравнением.}
Если
\[
\nu>0,
\]
то разрыв заменяется гладким переходным слоем конечной толщины. В пределе
\[
\nu\to 0+
\]
этот слой становится всё уже и приводит к ударной волне невязкого уравнения.

\paragraph{Мнемоническое правило / Суть на пальцах.}
Градиентная катастрофа --- это момент рождения разрыва. Ударная волна --- это уже «жизнь после катастрофы», когда волна не просто крутая, а действительно имеет скачок.

\section*{Практические задачи к разделу 3}

\subsection*{Задача 3(a). Невязкое уравнение Бюргерса при \texorpdfstring{$u(x,0)=\sin x$}{u(x,0)=sin x}}

\textbf{Дано.}
\[
u_t+u\,u_x=0,
\qquad
u(x,0)=\sin x.
\]

\textbf{Теоретическое обоснование.}
Это невязкое уравнение Бюргерса, то есть квазилинейное уравнение первого порядка. Для него применим метод характеристик. Вдоль характеристик решение постоянно, а сами характеристики зависят от начального значения решения.

\Solution
\textbf{Шаг 1. Записываем характеристическую систему.}

Для уравнения
\[
u_t+u\,u_x=0
\]
имеем
\[
\frac{dx}{dt}=u,
\qquad
\frac{du}{dt}=0.
\]

\textbf{Комментарий.} Второе уравнение означает, что вдоль характеристики величина $u$ сохраняется.

\textbf{Шаг 2. Параметризуем характеристики начальной координатой.}

Пусть характеристика выходит из точки
\[
x=\xi
\]
при
\[
t=0.
\]
Тогда по начальному условию
\[
u(\xi,0)=\sin \xi.
\]

Поскольку $u$ постоянно на характеристике, получаем
\[
u=\sin \xi.
\]

\textbf{Шаг 3. Находим уравнение характеристики.}

Подставляем найденное значение в уравнение для $x(t)$:
\[
\frac{dx}{dt}=\sin \xi.
\]

Интегрируя по $t$, получаем
\[
x=\xi+t\sin \xi.
\]

\textbf{Шаг 4. Записываем решение в параметрическом виде.}

Итак,
\[
u(x,t)=\sin \xi,
\qquad
x=\xi+t\sin \xi.
\]

Это и есть классическое решение задачи Коши в параметрическом представлении.

\textbf{Шаг 5. Получаем неявную форму решения.}

Из равенства
\[
u=\sin \xi
\]
и
\[
x=\xi+t\sin \xi=\xi+t u
\]
получаем
\[
\xi=x-tu.
\]

Подставляя это в формулу для $u$, имеем
\[
u(x,t)=\sin(x-tu(x,t)).
\]

Следовательно, решение может быть записано в \textbf{неявном виде}
\[
u=\sin(x-tu).
\]

\textbf{Шаг 6. Исследуем момент градиентной катастрофы.}

Общая формула даёт
\[
x_\xi=1+t\varphi'(\xi).
\]
Здесь
\[
\varphi(\xi)=\sin \xi,
\qquad
\varphi'(\xi)=\cos \xi.
\]
Поэтому
\[
x_\xi=1+t\cos \xi.
\]

Классическое решение существует, пока
\[
1+t\cos \xi>0
\]
для всех $\xi$.

Поскольку
\[
\min_\xi \cos \xi=-1,
\]
первый момент разрушения гладкости определяется условием
\[
1-t=0.
\]
Отсюда
\[
t_*=1.
\]

\textbf{Шаг 7. Укажем точку возникновения катастрофы.}

Равенство
\[
\cos \xi=-1
\]
выполняется при
\[
\xi=\pi+2\pi k,
\qquad k\in\mathbb{Z}.
\]
Для этих значений
\[
u=\sin \xi=0,
\]
поэтому
\[
x=\xi+t\sin \xi=\xi.
\]

Следовательно, при $t=1$ градиентная катастрофа возникает в точках
\[
x=\pi+2\pi k,
\qquad k\in\mathbb{Z}.
\]

\textbf{Ответ.}
\[
u(x,t)=\sin \xi,
\qquad
x=\xi+t\sin \xi,
\]
или, эквивалентно, в неявной форме
\[
u=\sin(x-tu).
\]
Это классическое решение существует при
\[
0\le t<1.
\]
Первый момент градиентной катастрофы:
\[
t_*=1.
\]

\subsection*{Задача 3(b). Почему при \texorpdfstring{$u_t+u u_x=\nu u_{xx}$}{ut+u ux=nu uxx} катастрофы нет}

\textbf{Дано.}
\[
u_t+u\,u_x=\nu u_{xx},
\qquad \nu>0.
\]

\textbf{Нужно объяснить.} Почему градиентная катастрофа не происходит и как ведёт себя решение качественно.

\textbf{Теоретическое обоснование.}
В невязком случае образование катастрофы связано с пересечением характеристик. В вязком случае появляется параболический диффузионный член $\nu u_{xx}$, который постоянно сглаживает решение и не даёт производной взорваться так, как это происходит в чисто транспортной модели.

\Solution
\textbf{Шаг 1. Сравним с невязким случаем.}

Для уравнения
\[
u_t+u\,u_x=0
\]
более высокие участки профиля движутся быстрее, и фронт крутизны растёт без ограничений. Это приводит к условию пересечения характеристик и к градиентной катастрофе.

\textbf{Шаг 2. Учитываем диффузионное слагаемое.}

В вязком уравнении присутствует член
\[
\nu u_{xx}.
\]

\textbf{Комментарий.} Это слагаемое описывает диффузию, то есть стремление профиля к выравниванию.

\textbf{Шаг 3. Объясняем механизм сглаживания.}

Если в некоторой точке градиент становится слишком большим, то вторая производная по $x$ тоже становится большой по модулю. Тогда член
\[
\nu u_{xx}
\]
начинает существенно влиять на эволюцию и сглаживает резкий фронт. Поэтому нелинейный перенос не успевает создать вертикальный обрыв.

\textbf{Шаг 4. Качественное поведение решения.}

Решение вязкого уравнения Бюргерса:
\begin{itemize}
  \item остаётся гладким;
  \item резкие перепады превращает в переходные слои конечной толщины;
  \item со временем сглаживает мелкие структуры;
  \item при больших временах стремится к более ровному распределению, если нет внешней подкачки.
\end{itemize}

\textbf{Шаг 5. Связь с ударной волной.}

Если сравнивать с невязким случаем, то там после градиентной катастрофы возникает разрыв. При $\nu>0$ вместо разрыва возникает \textbf{гладкий, но крутой фронт}. Его толщина тем меньше, чем меньше $\nu$, но при любом фиксированном $\nu>0$ этот фронт остаётся непрерывным и гладким.

\textbf{Ответ.}
Градиентная катастрофа не происходит, потому что диффузионный член
\[
\nu u_{xx}
\]
сглаживает растущие градиенты и препятствует пересечению характеристического механизма в чистом виде. Качественно решение образует не разрыв, а гладкий переходный слой конечной толщины, причём вязкость рассеивает резкие перепады.

\section*{4. Метод Фурье для параболического уравнения (теплопроводности)}

\textit{Замечание по источнику. Раздел составлен по гл.~5, \S\S~5.1--5.2 пособия Г.\,В.~Алексеева. Сохраняем принятые там обозначения: длина отрезка равна $l$, коэффициент температуропроводности обозначается через $a^2>0$, а цилиндр области записывается в виде $Q_T=(0,l)\times(0,T]$.}

\subsection*{4(a). Первая краевая задача для одномерного уравнения теплопроводности}

\Answer
\textbf{Первой краевой задачей} для одномерного уравнения теплопроводности на отрезке $(0,l)$ называется задача нахождения функции
\[
u\in C^{2,1}(Q_T)\cap C(\overline Q_T),
\qquad
Q_T=(0,l)\times(0,T],
\]
удовлетворяющей уравнению
\[
\frac{\partial u}{\partial t}=a^2\frac{\partial^2u}{\partial x^2}+f(x,t)
\qquad \text{в } Q_T,
\]
граничным условиям
\[
u(0,t)=g_1(t),
\qquad
u(l,t)=g_2(t),
\qquad
t\in(0,T],
\]
и начальному условию
\[
u(x,0)=\varphi(x),
\qquad
x\in(0,l).
\]

\textbf{Смысл входящих данных.}
\begin{itemize}
  \item $u(x,t)$ --- искомая температура;
  \item $a^2>0$ --- коэффициент температуропроводности;
  \item $f(x,t)$ --- плотность внутренних источников тепла;
  \item $g_1(t)$ и $g_2(t)$ --- температуры на концах стержня;
  \item $\varphi(x)$ --- начальное распределение температуры.
\end{itemize}

\textbf{Почему задача называется первой краевой.}
На границе задаются \textbf{значения самой функции} $u$, то есть условия Дирихле. Именно это и соответствует первой краевой задаче в терминологии Алексеева.

\textbf{Условие согласования в угловых точках.}
Если мы хотим искать \textbf{классическое} решение, то в точках пересечения начальной линии с боковой границей должны выполняться условия
\[
\varphi(0)=g_1(0),
\qquad
\varphi(l)=g_2(0).
\]
Иначе в точках $(0,0)$ и $(l,0)$ решение не сможет быть непрерывным одновременно и по начальному, и по граничному данным.

\textbf{Однородный частный случай.}
Для применения метода Фурье прежде всего рассматривают более простую задачу
\[
\frac{\partial u}{\partial t}=a^2\frac{\partial^2u}{\partial x^2},
\qquad
u(0,t)=u(l,t)=0,
\qquad
u(x,0)=\varphi(x).
\]
Именно к ней затем сводится общий случай после специальной замены неизвестной функции.

\textbf{Замечание о корректности.}
В \S~5.1 единственность и непрерывная зависимость решения от данных выводятся из \textbf{принципа максимума}. Поэтому для уравнения теплопроводности первая краевая задача имеет существенно более жёсткие свойства единственности, чем соответствующие гиперболические задачи.

\paragraph{Мнемоническое правило / Суть на пальцах.}
Первая краевая задача для теплопроводности --- это ситуация, когда нам заранее известны температуры на концах стержня и известна начальная температура внутри. Дальше уравнение описывает, как тепло «растекается» между этими фиксированными граничными значениями.

\subsection*{4(b). Как применяется метод разделения переменных и какая возникает спектральная задача}

\Answer
\textbf{Шаг 1. Сначала сводим общую задачу к задаче с однородными граничными условиями.}

Для общей задачи
\[
u_t=a^2u_{xx}+f(x,t),
\qquad
u(0,t)=g_1(t),
\qquad
u(l,t)=g_2(t),
\qquad
u(x,0)=\varphi(x)
\]
вводим функцию
\[
w(x,t)=g_1(t)+\bigl(g_2(t)-g_1(t)\bigr)\frac{x}{l}.
\]

\textbf{Пояснение.} Эта функция линейна по $x$, поэтому автоматически удовлетворяет граничным условиям
\[
w(0,t)=g_1(t),
\qquad
w(l,t)=g_2(t).
\]

Ищем решение в виде
\[
u(x,t)=v(x,t)+w(x,t).
\]
Тогда новая функция $v$ удовлетворяет однородным граничным условиям
\[
v(0,t)=u(0,t)-w(0,t)=0,
\qquad
v(l,t)=u(l,t)-w(l,t)=0.
\]
Начальное условие для $v$ имеет вид
\[
v(x,0)=\overline{\varphi}(x)=\varphi(x)-w(x,0).
\]

Подставим $u=v+w$ в исходное уравнение:
\[
v_t+w_t=a^2(v_{xx}+w_{xx})+f(x,t).
\]
Переносим все известные слагаемые вправо:
\[
v_t=a^2v_{xx}+\overline f(x,t),
\]
где
\[
\overline f(x,t)=f(x,t)-w_t+a^2w_{xx}.
\]
Так как $w$ линейна по $x$, то фактически
\[
w_{xx}=0,
\]
но мы сохраняем запись Алексеева в общем виде.

\textbf{Шаг 2. Выделяем основную однородную модель.}

После этой редукции центральной становится задача
\[
u_t=a^2u_{xx},
\qquad
u(0,t)=u(l,t)=0,
\qquad
u(x,0)=\varphi(x).
\]
Именно для неё применяется \textbf{метод разделения переменных}.

\textbf{Шаг 3. Ищем частные решения в виде произведения.}

Предположим, что решение имеет вид
\[
u(x,t)=X(x)T(t).
\]
Подставим это в уравнение:
\[
X(x)T'(t)=a^2X''(x)T(t).
\]
Перенесём множители так, чтобы левая часть зависела только от $t$, а правая --- только от $x$:
\[
\frac{T'(t)}{a^2T(t)}=\frac{X''(x)}{X(x)}.
\]

\textbf{Пояснение.} Левая часть зависит только от $t$, правая --- только от $x$. Значит, обе они должны быть равны одной и той же постоянной. По принятой в книге нормировке эту постоянную записывают в виде $-\lambda$:
\[
\frac{T'(t)}{a^2T(t)}=\frac{X''(x)}{X(x)}=-\lambda.
\]

Отсюда получаем два обыкновенных дифференциальных уравнения:
\[
T'(t)+a^2\lambda T(t)=0
\]
и
\[
X''(x)+\lambda X(x)=0.
\]

\textbf{Шаг 4. Переносим граничные условия на пространственный множитель.}

Так как
\[
u(0,t)=X(0)T(t)=0,
\qquad
u(l,t)=X(l)T(t)=0,
\]
то для нетривиального решения $T\not\equiv 0$ должны выполняться условия
\[
X(0)=0,
\qquad
X(l)=0.
\]

Следовательно, возникает \textbf{спектральная задача}
\[
X''+\lambda X=0
\quad \text{в } (0,l),
\qquad
X(0)=X(l)=0.
\]

\textbf{Почему здесь появляется спектр.}
Не всякое значение $\lambda$ допускает ненулевое решение краевой задачи для $X$. Допустимы только специальные значения $\lambda=\lambda_k$, называемые \textbf{собственными значениями}. Соответствующие функции $X_k$ называются \textbf{собственными функциями}. Именно по ним затем строится ряд Фурье для решения.

\paragraph{Мнемоническое правило / Суть на пальцах.}
Метод Фурье раскладывает сложный тепловой процесс на сумму простых «мод».
Каждая мода имеет фиксированную форму по $x$ и отдельно затухает по времени.
Пространственные формы разрешаются не произвольно, а только теми, которые совместимы с граничными условиями.

\subsection*{4(c). Собственные значения и собственные функции при условиях Дирихле}

\Answer
Нужно решить спектральную задачу
\[
X''+\lambda X=0,
\qquad
X(0)=X(l)=0.
\]
Разберём её подробно по знаку параметра $\lambda$.

\textbf{Случай 1: \texorpdfstring{$\lambda=0$}{lambda=0}.}

Тогда
\[
X''=0,
\]
откуда
\[
X(x)=C_1x+C_2.
\]
Из условия
\[
X(0)=0
\]
получаем
\[
C_2=0.
\]
Из условия
\[
X(l)=0
\]
получаем
\[
C_1l=0.
\]
Следовательно,
\[
C_1=0,
\qquad X\equiv 0.
\]
Нетривиального решения нет.

\textbf{Случай 2: \texorpdfstring{$\lambda<0$}{lambda<0}.}

Положим
\[
\lambda=-\mu^2,
\qquad
\mu>0.
\]
Тогда уравнение принимает вид
\[
X''-\mu^2X=0.
\]
Его общее решение:
\[
X(x)=C_1e^{\mu x}+C_2e^{-\mu x}
\]
или эквивалентно
\[
X(x)=A\cosh(\mu x)+B\sinh(\mu x).
\]
Из условия
\[
X(0)=0
\]
следует
\[
A=0.
\]
Тогда
\[
X(x)=B\sinh(\mu x).
\]
Но из
\[
X(l)=B\sinh(\mu l)=0
\]
и неравенства
\[
\sinh(\mu l)\ne 0
\]
получаем
\[
B=0.
\]
Снова остаётся только тривиальное решение.

\textbf{Случай 3: \texorpdfstring{$\lambda>0$}{lambda>0}.}

Положим
\[
\lambda=\mu^2,
\qquad
\mu>0.
\]
Тогда
\[
X''+\mu^2X=0,
\]
и общее решение имеет вид
\[
X(x)=C_1\cos(\mu x)+C_2\sin(\mu x).
\]

Из условия
\[
X(0)=0
\]
получаем
\[
C_1=0.
\]
Значит,
\[
X(x)=C_2\sin(\mu x).
\]

Теперь используем условие
\[
X(l)=0:
\qquad
C_2\sin(\mu l)=0.
\]
Для нетривиального решения $C_2\ne 0$, поэтому необходимо
\[
\sin(\mu l)=0.
\]
Это равносильно условию
\[
\mu l=k\pi,
\qquad
k=1,2,\dots
\]
Следовательно,
\[
\mu_k=\frac{k\pi}{l},
\qquad
\lambda_k=\mu_k^2=\left(\frac{k\pi}{l}\right)^2.
\]

Соответствующие собственные функции:
\[
X_k(x)=\sin\frac{k\pi x}{l},
\qquad
k=1,2,\dots
\]

\textbf{Итак, окончательный ответ для спектральной задачи:}
\[
\lambda_k=\left(\frac{k\pi}{l}\right)^2,
\qquad
X_k(x)=\sin\frac{k\pi x}{l},
\qquad
k=1,2,\dots
\]

\textbf{Ортогональность.}
Эти функции образуют ортогональную систему на $(0,l)$:
\[
\int_0^l
\sin\frac{k\pi x}{l}\,
\sin\frac{m\pi x}{l}\,dx
=0,
\qquad
k\ne m.
\]
Кроме того,
\[
\int_0^l \sin^2\frac{k\pi x}{l}\,dx=\frac{l}{2}.
\]
Именно эта ортогональность позволяет однозначно находить коэффициенты ряда Фурье.

\paragraph{Мнемоническое правило / Суть на пальцах.}
При нулевой температуре на концах стержня разрешены только синусоидальные формы, которые сами обращаются в нуль на концах. Чем больше номер $k$, тем более «мелкая» пространственная волна получается.

\subsection*{4(d). Вид решения в форме ряда Фурье и формулы для коэффициентов}

\Answer
\textbf{Шаг 1. Временная часть для каждой моды.}

После нахождения собственных значений
\[
\lambda_k=\left(\frac{k\pi}{l}\right)^2
\]
уравнение для временного множителя принимает вид
\[
T_k'(t)+a^2\lambda_k T_k(t)=0.
\]
Это линейное однородное уравнение первого порядка. Его решение:
\[
T_k(t)=a_ke^{-a^2\lambda_kt}
=a_k\exp\!\left[-\left(\frac{k\pi a}{l}\right)^2t\right].
\]
Здесь $a_k$ --- пока неизвестные постоянные.

\textbf{Шаг 2. Строим элементарные моды.}

Для каждого $k$ получаем частное решение
\[
u_k(x,t)=a_ke^{-a^2\lambda_kt}\sin\frac{k\pi x}{l}.
\]
По линейности уравнения можно брать сумму таких мод. Поэтому ищем решение в виде ряда
\[
u(x,t)=\sum_{k=1}^{\infty}
a_ke^{-a^2\lambda_kt}\sin\frac{k\pi x}{l}.
\]
Подставляя явный вид $\lambda_k$, получаем
\[
u(x,t)=\sum_{k=1}^{\infty}
a_k\exp\!\left[-\left(\frac{k\pi a}{l}\right)^2t\right]
\sin\frac{k\pi x}{l}.
\]

\textbf{Шаг 3. Определяем коэффициенты из начального условия.}

При $t=0$ должно выполняться
\[
u(x,0)=\varphi(x).
\]
Поэтому
\[
\varphi(x)=\sum_{k=1}^{\infty}a_k\sin\frac{k\pi x}{l}.
\]
Это разложение функции $\varphi$ в \textbf{ряд Фурье по синусам} на интервале $(0,l)$.

\textbf{Шаг 4. Выводим формулу для \texorpdfstring{$a_k$}{ak}.}

Умножим обе части разложения на
\[
\sin\frac{m\pi x}{l}
\]
и проинтегрируем по $x$ от $0$ до $l$:
\[
\int_0^l \varphi(x)\sin\frac{m\pi x}{l}\,dx
=
\sum_{k=1}^{\infty}
a_k
\int_0^l
\sin\frac{k\pi x}{l}\sin\frac{m\pi x}{l}\,dx.
\]
По ортогональности все слагаемые с $k\ne m$ исчезают, и остаётся
\[
\int_0^l \varphi(x)\sin\frac{m\pi x}{l}\,dx
=
a_m\int_0^l \sin^2\frac{m\pi x}{l}\,dx
=a_m\frac{l}{2}.
\]
Отсюда
\[
a_m=\frac{2}{l}\int_0^l \varphi(x)\sin\frac{m\pi x}{l}\,dx.
\]
Переименовывая индекс $m$ обратно в $k$, получаем формулу Алексеева:
\[
a_k=\frac{2}{l}\int_0^l \varphi(x)\sin\frac{k\pi x}{l}\,dx,
\qquad
k=1,2,\dots
\]

\textbf{Итак, решение первой однородной краевой задачи имеет вид}
\[
u(x,t)=\sum_{k=1}^{\infty}
\left(
\frac{2}{l}\int_0^l \varphi(\xi)\sin\frac{k\pi \xi}{l}\,d\xi
\right)
\exp\!\left[-\left(\frac{k\pi a}{l}\right)^2t\right]
\sin\frac{k\pi x}{l}.
\]

\textbf{Условия, при которых этот ряд даёт классическое решение.}
В \S~5.2 используется условие
\[
\varphi\in C[0,l],
\qquad
\varphi' \text{ кусочно-непрерывна на } [0,l],
\qquad
\varphi(0)=\varphi(l)=0.
\]
Тогда ряд для $u$ сходится равномерно, а при каждом $t>0$ его можно дифференцировать по $x$ и $t$ столько раз, сколько требуется для обоснования уравнения.

\textbf{Полезное дополнение для неоднородного уравнения.}
Если
\[
u_t=a^2u_{xx}+f(x,t),
\qquad
u(0,t)=u(l,t)=0,
\qquad
u(x,0)=0,
\]
то, по Алексееву, решение ищется в виде
\[
u(x,t)=\sum_{k=1}^{\infty}T_k(t)\sin\frac{k\pi x}{l},
\]
где правую часть тоже раскладывают в ряд Фурье:
\[
f(x,t)=\sum_{k=1}^{\infty}f_k(t)\sin\frac{k\pi x}{l},
\qquad
f_k(t)=\frac{2}{l}\int_0^l f(x,t)\sin\frac{k\pi x}{l}\,dx.
\]
Тогда
\[
T_k'(t)+\omega_k^2T_k(t)=f_k(t),
\qquad
\omega_k=\frac{k\pi a}{l},
\qquad
T_k(0)=0,
\]
и
\[
T_k(t)=\int_0^t e^{-\omega_k^2(t-\tau)}f_k(\tau)\,d\tau.
\]
Следовательно,
\[
u(x,t)=\sum_{k=1}^{\infty}
\left[
\int_0^t e^{-\omega_k^2(t-\tau)}f_k(\tau)\,d\tau
\right]
\sin\frac{k\pi x}{l}.
\]

\paragraph{Мнемоническое правило / Суть на пальцах.}
Ряд Фурье для теплопроводности --- это разложение начальной температуры по синусам. Каждый синус живёт сам по себе и убывает по экспоненте. Коэффициенты Фурье просто измеряют, «сколько» каждой моды было в начальном профиле.

\subsection*{4(e). Почему для параболического уравнения нужно одно начальное условие, а для гиперболического --- два}

\Answer
\textbf{Ключевое различие состоит в порядке уравнения по времени.}

Для уравнения теплопроводности
\[
u_t=a^2u_{xx}
\]
производная по времени имеет \textbf{первый порядок}. После разделения переменных для каждой моды получается уравнение
\[
T_k'(t)+a^2\lambda_kT_k(t)=0.
\]
Это обыкновенное дифференциальное уравнение \textbf{первого порядка}, а значит, его общее решение содержит \textbf{одну} произвольную постоянную:
\[
T_k(t)=a_ke^{-a^2\lambda_kt}.
\]
Чтобы определить эту постоянную, достаточно одного условия при $t=0$, то есть знания
\[
u(x,0)=\varphi(x).
\]

Для волнового уравнения
\[
u_{tt}=a^2u_{xx}
\]
после разделения переменных возникает уравнение
\[
T_k''(t)+a^2\lambda_kT_k(t)=0,
\]
то есть уравнение \textbf{второго порядка} по времени. Его общее решение уже содержит две произвольные постоянные. Поэтому требуется два начальных условия:
\[
u(x,0)=\varphi(x),
\qquad
u_t(x,0)=\psi(x).
\]

\textbf{Физическая интерпретация.}
\begin{itemize}
  \item Для теплопроводности достаточно знать текущее распределение температуры: дальше эволюция определяется законом диффузии.
  \item Для колебаний этого мало: нужно знать не только положение, но и начальную скорость, иначе движение определить нельзя.
\end{itemize}

\textbf{Связь с поведением решений.}
Это различие не формально. Уравнение теплопроводности описывает \textbf{затухающий, сглаживающий} процесс, а волновое уравнение --- \textbf{инерционный} процесс с памятью о начальной скорости.

\paragraph{Мнемоническое правило / Суть на пальцах.}
Тепло не «разгоняется», а просто перераспределяется, поэтому нужно знать только, \textit{как нагрето сейчас}. Волна же ещё и летит по инерции, поэтому нужно знать не только положение, но и \textit{как быстро всё стартовало}.

\subsection*{4(f). Поведение решения при \texorpdfstring{$t\to\infty$}{t->infinity}: затухание мод}

\Answer
Пусть рассматривается однородная первая краевая задача
\[
u_t=a^2u_{xx},
\qquad
u(0,t)=u(l,t)=0,
\qquad
u(x,0)=\varphi(x).
\]
Её решение имеет вид
\[
u(x,t)=\sum_{k=1}^{\infty}
a_k\exp\!\left[-\left(\frac{k\pi a}{l}\right)^2t\right]
\sin\frac{k\pi x}{l}.
\]

\textbf{Шаг 1. Рассмотрим поведение одной моды.}

Для каждого фиксированного $k$ множитель
\[
\exp\!\left[-\left(\frac{k\pi a}{l}\right)^2t\right]
\]
стремится к нулю при
\[
t\to\infty.
\]
Следовательно, каждая отдельная гармоника затухает.

\textbf{Шаг 2. Сравним скорости затухания разных мод.}

Показатель экспоненты содержит множитель
\[
\left(\frac{k\pi a}{l}\right)^2.
\]
Чем больше $k$, тем больше этот множитель, а значит, тем быстрее убывает соответствующая мода.

\textbf{Вывод.} Высокочастотные, мелкомасштабные пространственные колебания исчезают быстрее, чем низкочастотные.

\textbf{Шаг 3. Поведение всего решения.}

При однородных граничных условиях и отсутствии источников все моды затухают, поэтому
\[
u(x,t)\to 0
\qquad \text{при } t\to\infty.
\]
Физически это означает, что стержень со временем остывает до нулевой температуры на концах и во всём объёме.

\textbf{Если задача неоднородна.}
Если присутствуют ненулевые граничные условия или источники, то после выделения стационарной части решение представляется как
\[
u=v+w,
\]
где переходная часть $v$ раскладывается по модам и затухает, а предельное поведение определяется функцией $w$, то есть стационарным тепловым режимом.

\paragraph{Мнемоническое правило / Суть на пальцах.}
Теплопроводность «убивает рябь».
Сначала быстро исчезают мелкие неровности температуры,
потом медленно затухают крупные.
В конце остаётся либо ноль, либо стационарный профиль,
если его поддерживает граница или источник.

\section*{Практические задачи к разделу 4}

\subsection*{Задача 4(a). Решить \texorpdfstring{$u_t=u_{xx}$}{ut=uxx} на \texorpdfstring{$(0,\pi)$}{(0,pi)} при \texorpdfstring{$u(0,t)=u(\pi,t)=0$}{u(0,t)=u(pi,t)=0}, \texorpdfstring{$u(x,0)=\sin 2x$}{u(x,0)=sin 2x}}

\textbf{Дано.}
\[
u_t=u_{xx},
\qquad
0<x<\pi,
\qquad
t>0,
\]
\[
u(0,t)=u(\pi,t)=0,
\qquad
u(x,0)=\sin 2x.
\]

\textbf{Теоретическое обоснование.}
Это первая однородная краевая задача для уравнения теплопроводности на отрезке $(0,\pi)$ при
\[
a=1,
\qquad
l=\pi.
\]
По методу Фурье решение ищется в виде ряда
\[
u(x,t)=\sum_{k=1}^{\infty}a_ke^{-k^2t}\sin kx,
\]
поскольку для $l=\pi$ имеем
\[
\lambda_k=\left(\frac{k\pi}{\pi}\right)^2=k^2,
\qquad
X_k(x)=\sin kx.
\]

\Solution
\textbf{Шаг 1. Запишем общее решение в виде ряда Фурье.}

Для нашей задачи
\[
u(x,t)=\sum_{k=1}^{\infty}a_ke^{-k^2t}\sin kx.
\]

\textbf{Шаг 2. Используем начальное условие.}

При $t=0$ должно выполняться
\[
\sin 2x=\sum_{k=1}^{\infty}a_k\sin kx.
\]

\textbf{Комментарий.} Начальная функция уже является одной собственной функцией спектральной задачи, а именно гармоникой с номером $k=2$.

\textbf{Шаг 3. Определяем коэффициенты Фурье.}

Из разложения сразу видно, что
\[
a_2=1,
\qquad
a_k=0 \quad \text{при } k\ne 2.
\]

\textbf{Шаг 4. Подставляем коэффициенты в общий ряд.}

Остаётся только одно слагаемое:
\[
u(x,t)=e^{-4t}\sin 2x.
\]

\textbf{Проверка.}

Проверим уравнение:
\[
u_t=-4e^{-4t}\sin 2x,
\]
\[
u_{xx}=-4e^{-4t}\sin 2x.
\]
Следовательно,
\[
u_t=u_{xx}.
\]

Проверим граничные условия:
\[
u(0,t)=e^{-4t}\sin 0=0,
\qquad
u(\pi,t)=e^{-4t}\sin 2\pi=0.
\]

Проверим начальное условие:
\[
u(x,0)=e^0\sin 2x=\sin 2x.
\]

\textbf{Ответ.}
\[
u(x,t)=e^{-4t}\sin 2x.
\]

\subsection*{Задача 4(b). Для \texorpdfstring{$u(x,0)=x(\pi-x)$}{u(x,0)=x(pi-x)} найти коэффициенты Фурье}

\textbf{Дано.}
\[
u_t=u_{xx},
\qquad
0<x<\pi,
\qquad
t>0,
\]
\[
u(0,t)=u(\pi,t)=0,
\qquad
u(x,0)=x(\pi-x).
\]

\textbf{Теоретическое обоснование.}
Это та же первая однородная краевая задача на $(0,\pi)$, поэтому решение имеет вид
\[
u(x,t)=\sum_{k=1}^{\infty}b_ke^{-k^2t}\sin kx.
\]
Нужно только правильно поставить формулы для коэффициентов $b_k$.

\Solution
\textbf{Шаг 1. Записываем разложение начальной функции по синусам.}

Из условия при $t=0$ получаем
\[
x(\pi-x)=\sum_{k=1}^{\infty}b_k\sin kx.
\]

\textbf{Шаг 2. Используем общую формулу коэффициентов для интервала \texorpdfstring{$(0,\pi)$}{(0,pi)}.}

Поскольку здесь
\[
l=\pi,
\]
формула
\[
a_k=\frac{2}{l}\int_0^l \varphi(x)\sin\frac{k\pi x}{l}\,dx
\]
переходит в
\[
b_k=\frac{2}{\pi}\int_0^{\pi}x(\pi-x)\sin kx\,dx,
\qquad
k=1,2,\dots
\]

\textbf{Шаг 3. Записываем решение, не вычисляя интегралы.}

Тогда
\[
u(x,t)=\sum_{k=1}^{\infty}
\left(
\frac{2}{\pi}\int_0^{\pi}x(\pi-x)\sin k\xi\,d\xi
\right)e^{-k^2t}\sin kx.
\]

\textbf{Комментарий.} По условию задачи вычислять интегралы не требуется; важно правильно поставить разложение и формулы для коэффициентов.

\textbf{Ответ.}
\[
b_k=\frac{2}{\pi}\int_0^{\pi}x(\pi-x)\sin kx\,dx,
\qquad
k=1,2,\dots
\]
и
\[
u(x,t)=\sum_{k=1}^{\infty}b_ke^{-k^2t}\sin kx.
\]

\subsection*{Задача 4(c). Почему при \texorpdfstring{$t>0$}{t>0} решение быстро становится гладким}

\textbf{Дано.}
Первая краевая задача для уравнения теплопроводности на отрезке с начальной функцией $\varphi$, которая может быть недостаточно гладкой.

\textbf{Нужно объяснить.}
Почему при любом фиксированном $t>0$ решение оказывается гораздо более гладким, чем начальная функция.

\textbf{Теоретическое обоснование.}
По методу Фурье решение представляется рядом
\[
u(x,t)=\sum_{k=1}^{\infty}
a_ke^{-\omega_k^2t}X_k(x),
\qquad
\omega_k=\frac{k\pi a}{l},
\]
где $X_k$ --- собственные функции спектральной задачи. Ключевой объект здесь --- экспоненциальный множитель
\[
e^{-\omega_k^2t}.
\]

\Solution
\textbf{Шаг 1. Смотрим, что происходит с высокими модами.}

Если номер моды $k$ велик, то число
\[
\omega_k^2=\left(\frac{k\pi a}{l}\right)^2
\]
тоже велико. Поэтому при любом фиксированном $t>0$
\[
e^{-\omega_k^2t}
\]
становится очень маленьким.

\textbf{Комментарий.} Именно высокие моды отвечают за резкие изломы, мелкие осцилляции и плохую гладкость начальной функции.

\textbf{Шаг 2. Проверяем производные ряда.}

При дифференцировании по $x$ каждая производная даёт множители порядка $k$, а при дифференцировании по $t$ --- множители порядка $k^2$. Поэтому типичный член продифференцированного ряда имеет вид
\[
k^m e^{-ck^2t}
\]
с некоторыми $m\in\mathbb{N}$ и $c>0$.

\textbf{Шаг 3. Сравниваем степенной рост и экспоненциальное затухание.}

Для любого фиксированного $t>0$ экспонента
\[
e^{-ck^2t}
\]
убывает быстрее любой степени $k^m$. Значит, ряды, полученные дифференцированием решения по $x$ и по $t$, сходятся абсолютно и равномерно.

\textbf{Шаг 4. Делаем вывод о гладкости.}

Раз продифференцированные ряды сходятся, то решение можно дифференцировать сколько угодно раз при любом $t>0$. Следовательно, даже если начальная функция имела изломы или только конечное число производных, решение мгновенно становится очень гладким внутри области.

\textbf{Шаг 5. Физический смысл.}

Теплопроводность немедленно сглаживает резкие перепады температуры. Поэтому «углы» начального профиля не сохраняются: они моментально размываются.

\textbf{Ответ.}
Потому что в ряде Фурье для решения каждая мода умножается на экспоненту
\[
e^{-\omega_k^2t},
\]
которая при любом $t>0$ резко подавляет высокочастотные гармоники. Из-за этого ряды для производных тоже сходятся, и решение становится бесконечно гладким по пространству и времени при любом положительном времени.

\section*{5. Метод Фурье для внутренней задачи уравнения Лапласа}

\textit{Замечание по источнику. Раздел составлен по гл.~6, \S~6.4 пособия Г.\,В.~Алексеева. В книге одновременно рассматриваются внутренняя и внешняя задачи Дирихле для круга. Ниже основной акцент сделан на внутренней задаче, поскольку именно она вынесена в список вопросов. Сохраняем обозначения Алексеева: радиус круга равен $a$, полярные координаты обозначаются через $(\rho,\varphi)$.}

\subsection*{5(a). Задача Дирихле для уравнения Лапласа в круге}

\Answer
\textbf{Постановка области.}
Рассматривается круг радиуса $a$:
\[
\Omega=\Omega_a=\{(x,y)\in\mathbb{R}^2:\ x^2+y^2<a^2\},
\qquad
\Gamma_a=\partial\Omega.
\]

\textbf{Внутренняя задача Дирихле} для уравнения Лапласа состоит в нахождении классического решения
\[
u\in C^2(\Omega)\cap C(\overline\Omega)
\]
уравнения
\[
\Delta u=0
\qquad \text{в } \Omega,
\]
удовлетворяющего граничному условию
\[
u=g
\qquad \text{на } \Gamma_a.
\]

\textbf{Что означает граничное условие на окружности.}
Если на границе перейти к угловой параметризации
\[
x=a\cos\varphi,
\qquad
y=a\sin\varphi,
\qquad
\varphi\in[0,2\pi],
\]
то граничное условие записывается как
\[
u(a,\varphi)=g(\varphi).
\]

\textbf{Требования к граничной функции.}
Для корректной постановки граничная функция должна быть $2\pi$-периодической:
\[
g(0)=g(2\pi).
\]
В \S~6.4 сначала предполагается
\[
g\in C[0,2\pi],
\qquad
g(0)=g(2\pi),
\]
а для перехода к классическому решению дополнительно используется условие кусочной непрерывности производной $g'$.

\textbf{Почему это именно задача Дирихле.}
На границе задаётся \textbf{значение самой функции} $u$, а не нормальной производной. Это и есть первая краевая задача для эллиптического уравнения.

\textbf{Замечание о единственности.}
Единственность решения внутренней задачи Дирихле для гармонических функций следует из \textbf{принципа максимума}: если две гармонические функции совпадают на границе, то их разность гармонична и равна нулю на границе, а значит, равна нулю во всём круге.

\paragraph{Мнемоническое правило / Суть на пальцах.}
Внутренняя задача Дирихле в круге отвечает на вопрос: если по окружности задана температура или потенциал, то каким будет стационарное поле внутри круга? Граница полностью «диктует» внутреннее гармоническое состояние.

\subsection*{5(b). Метод разделения переменных в полярных координатах. Радиальная и угловая части}

\Answer
\textbf{Шаг 1. Переходим к полярным координатам.}

На плоскости вводим координаты
\[
x=\rho\cos\varphi,
\qquad
y=\rho\sin\varphi,
\qquad
0\le \rho<a,
\qquad
0\le \varphi\le 2\pi.
\]
В этих переменных оператор Лапласа принимает вид
\[
\Delta u
=
\frac{1}{\rho}\frac{\partial}{\partial\rho}
\left(
\rho\frac{\partial u}{\partial\rho}
\right)
\frac{1}{\rho^2}\frac{\partial^2u}{\partial\varphi^2}.
\]
Поэтому уравнение Лапласа записывается так:
\[
\frac{1}{\rho}\frac{\partial}{\partial\rho}
\left(
\rho\frac{\partial u}{\partial\rho}
\right)
\frac{1}{\rho^2}\frac{\partial^2u}{\partial\varphi^2}=0.
\]

\textbf{Шаг 2. Ищем частные решения в виде произведения.}

Положим
\[
u(\rho,\varphi)=R(\rho)\Phi(\varphi).
\]
Подставим это выражение в уравнение:
\[
\frac{1}{\rho}\frac{d}{d\rho}\bigl(\rho R'(\rho)\Phi(\varphi)\bigr)
\frac{1}{\rho^2}R(\rho)\Phi''(\varphi)=0.
\]

Так как $\Phi$ не зависит от $\rho$, а $R$ не зависит от $\varphi$, имеем
\[
\frac{1}{\rho}\Phi(\varphi)(\rho R'(\rho))'
\frac{1}{\rho^2}R(\rho)\Phi''(\varphi)=0.
\]
Умножим на $\rho^2$:
\[
\rho\,\Phi(\varphi)(\rho R'(\rho))'
R(\rho)\Phi''(\varphi)=0.
\]
Теперь разделим на $R(\rho)\Phi(\varphi)$:
\[
\frac{\rho(\rho R')'}{R}+\frac{\Phi''}{\Phi}=0.
\]
Перенесём второе слагаемое вправо:
\[
\frac{\rho(\rho R')'}{R}
=
-\frac{\Phi''}{\Phi}.
\]

\textbf{Пояснение.} Левая часть зависит только от $\rho$, правая --- только от $\varphi$. Значит, обе части должны быть равны одной и той же постоянной. В обозначениях Алексеева:
\[
\frac{\rho(\rho R')'}{R}
=
-\frac{\Phi''}{\Phi}
=\lambda.
\]

Отсюда получаем два обыкновенных уравнения:
\[
\rho(\rho R')'-\lambda R=0
\qquad \text{в } (0,a),
\]
\[
\Phi''+\lambda\Phi=0
\qquad \text{в } (0,2\pi).
\]

\textbf{Шаг 3. Учитываем периодичность по углу.}

Поскольку функция $u(\rho,\varphi)$ должна быть однозначной на плоскости, значение при углах $\varphi=0$ и $\varphi=2\pi$ должно совпадать. Поэтому
\[
\Phi(0)=\Phi(2\pi).
\]
Фактически здесь также подразумевается согласованность производных, и именно это приводит к стандартной $2\pi$-периодической спектральной задаче для угловой части.

\textbf{Итак, после разделения переменных получаются:}
\begin{itemize}
  \item \textbf{угловое уравнение}
  \[
  \Phi''+\lambda\Phi=0,
  \qquad
  \Phi(0)=\Phi(2\pi),
  \]
  \item \textbf{радиальное уравнение}
  \[
  \rho(\rho R')'-\lambda R=0.
  \]
\end{itemize}

\textbf{Роль этих двух уравнений.}
Угловая задача определяет допустимые \textbf{спектральные параметры} $\lambda$, а затем для каждого такого $\lambda$ решается радиальное уравнение. После этого линейной комбинацией найденных мод подгоняют граничное условие $u(a,\varphi)=g(\varphi)$.

\paragraph{Мнемоническое правило / Суть на пальцах.}
В круге естественно отделить расстояние до центра и направление. Угловая часть отвечает за «рисунок по окружности», а радиальная --- за то, как этот рисунок протягивается от границы к центру.

\subsection*{5(c). Собственные значения и собственные функции угловой части}

\Answer
Рассмотрим спектральную задачу
\[
\Phi''+\lambda\Phi=0,
\qquad
\Phi(0)=\Phi(2\pi).
\]
Разберём её по знаку параметра $\lambda$.

\textbf{Случай 1: \texorpdfstring{$\lambda=0$}{lambda=0}.}

Тогда
\[
\Phi''=0,
\]
и общее решение имеет вид
\[
\Phi(\varphi)=A+B\varphi.
\]
Чтобы функция была $2\pi$-периодической, необходимо
\[
\Phi(0)=\Phi(2\pi).
\]
Подставляя, получаем
\[
A=A+2\pi B,
\]
откуда
\[
B=0.
\]
Следовательно, при $\lambda=0$ остаются только постоянные функции:
\[
\Phi_0(\varphi)=A.
\]

\textbf{Случай 2: \texorpdfstring{$\lambda<0$}{lambda<0}.}

Положим
\[
\lambda=-\mu^2,
\qquad \mu>0.
\]
Тогда уравнение принимает вид
\[
\Phi''-\mu^2\Phi=0,
\]
и его общее решение записывается как
\[
\Phi(\varphi)=A e^{\mu\varphi}+B e^{-\mu\varphi}.
\]
Нетривиальная $2\pi$-периодичность для такой функции невозможна, поскольку экспоненты не возвращаются к исходному значению при сдвиге на $2\pi$. Значит, в этом случае нет нетривиальных собственных функций.

\textbf{Случай 3: \texorpdfstring{$\lambda>0$}{lambda>0}.}

Положим
\[
\lambda=\mu^2,
\qquad \mu>0.
\]
Тогда
\[
\Phi''+\mu^2\Phi=0,
\]
и общее решение:
\[
\Phi(\varphi)=A\cos\mu\varphi+B\sin\mu\varphi.
\]
Чтобы функция была $2\pi$-периодической, нужно, чтобы
\[
\cos\mu(\varphi+2\pi)=\cos\mu\varphi,
\qquad
\sin\mu(\varphi+2\pi)=\sin\mu\varphi.
\]
Это выполняется тогда и только тогда, когда
\[
\mu=k,
\qquad
k=1,2,\dots
\]
Следовательно,
\[
\lambda_k=k^2,
\qquad
k=1,2,\dots
\]
а соответствующие собственные функции имеют вид
\[
\Phi_k(\varphi)=a_k\cos k\varphi+b_k\sin k\varphi.
\]

\textbf{Итог.}
Спектральная задача для угловой части имеет собственные значения
\[
\lambda_k=k^2,
\qquad
k=0,1,2,\dots
\]
а собственные функции:
\[
\Phi_0(\varphi)=\text{const},
\]
\[
\Phi_k(\varphi)=a_k\cos k\varphi+b_k\sin k\varphi,
\qquad
k\ge 1.
\]

\textbf{Почему именно эти функции важны.}
Они образуют тригонометрическую систему на окружности, а значит, позволяют разложить граничную функцию $g(\varphi)$ в ряд Фурье. Именно это и делает возможным согласование решения с граничным условием.

\paragraph{Мнемоническое правило / Суть на пальцах.}
Угол живёт по окружности, поэтому его собственные функции обязаны быть периодическими. Отсюда автоматически возникают косинусы и синусы --- единственные естественные гармоники на круговой границе.

\subsection*{5(d). Общее решение в виде ряда по косинусам и синусам}

\Answer
\textbf{Шаг 1. Решаем радиальное уравнение для каждого \texorpdfstring{$k$}{k}.}

После подстановки
\[
\lambda_k=k^2
\]
в радиальное уравнение получаем
\[
\rho(\rho R')'-k^2R=0.
\]
Раскрывая производную, имеем
\[
\rho^2R''+\rho R'-k^2R=0.
\]
Это уравнение Эйлера. Будем искать решение в виде степени
\[
R(\rho)=\rho^{\mu}.
\]
Тогда
\[
R'(\rho)=\mu \rho^{\mu-1},
\qquad
R''(\rho)=\mu(\mu-1)\rho^{\mu-2}.
\]
Подставим в уравнение:
\[
\rho^2\mu(\mu-1)\rho^{\mu-2}
\rho\mu\rho^{\mu-1}
-k^2\rho^\mu=0.
\]
Сокращая общий множитель $\rho^\mu$, получаем
\[
\mu(\mu-1)+\mu-k^2=0,
\]
то есть
\[
\mu^2-k^2=0.
\]
Следовательно,
\[
\mu=\pm k.
\]

\textbf{Итак, при \texorpdfstring{$k\ge 1$}{k>=1} имеем два линейно независимых радиальных решения:}
\[
\rho^k,
\qquad
\rho^{-k}.
\]
При
\[
k=0
\]
независимыми решениями являются
\[
1,
\qquad
\ln\rho.
\]

\textbf{Шаг 2. Выбираем те решения, которые допустимы внутри круга.}

Для \textbf{внутренней} задачи решение должно оставаться ограниченным при
\[
\rho\to 0.
\]
Поэтому:
\begin{itemize}
  \item функция $\rho^{-k}$ при $k\ge 1$ не подходит, так как она неограничена в центре;
  \item функция $\ln\rho$ тоже не подходит, так как при $\rho\to 0$ стремится к $-\infty$.
\end{itemize}
Значит, для внутренней задачи остаются только радиальные множители
\[
R_0(\rho)=1,
\qquad
R_k(\rho)=\rho^k,
\qquad
k\ge 1.
\]

\textbf{Шаг 3. Строим частные гармоники.}

Соединяя радиальную и угловую части, получаем частные решения
\[
u_k(\rho,\varphi)=\rho^k\bigl(a_k\cos k\varphi+b_k\sin k\varphi\bigr),
\qquad
k=0,1,2,\dots
\]
По линейности уравнения Лапласа сумма таких решений снова будет решением. Поэтому формально получаем ряд
\[
u(\rho,\varphi)=\sum_{k=0}^{\infty}
\rho^k\bigl(a_k\cos k\varphi+b_k\sin k\varphi\bigr).
\]

\textbf{Шаг 4. Подгоняем коэффициенты под граничное условие.}

На границе круга $\rho=a$ должно выполняться
\[
u(a,\varphi)=g(\varphi).
\]
Подставим $\rho=a$ в ряд:
\[
g(\varphi)=\sum_{k=0}^{\infty}
a^k\bigl(a_k\cos k\varphi+b_k\sin k\varphi\bigr).
\]

Разложим граничную функцию в ряд Фурье:
\[
g(\varphi)=\frac{\alpha_0}{2}+
\sum_{k=1}^{\infty}
\bigl(\alpha_k\cos k\varphi+\beta_k\sin k\varphi\bigr),
\]
где
\[
\alpha_0=\frac{1}{\pi}\int_0^{2\pi}g(\psi)\,d\psi,
\]
\[
\alpha_k=\frac{1}{\pi}\int_0^{2\pi}g(\psi)\cos k\psi\,d\psi,
\qquad
\beta_k=\frac{1}{\pi}\int_0^{2\pi}g(\psi)\sin k\psi\,d\psi,
\qquad
k=1,2,\dots
\]

Сравнивая коэффициенты, получаем
\[
a_0=\frac{\alpha_0}{2},
\qquad
a_k=\frac{\alpha_k}{a^k},
\qquad
b_k=\frac{\beta_k}{a^k},
\qquad
k\ge 1.
\]

\textbf{Следовательно, решение внутренней задачи Дирихле в круге имеет вид}
\[
u(\rho,\varphi)
=
\frac{\alpha_0}{2}
\sum_{k=1}^{\infty}
\left(\frac{\rho}{a}\right)^k
\bigl(\alpha_k\cos k\varphi+\beta_k\sin k\varphi\bigr).
\]
Это именно формула (4.23) у Алексеева.

\textbf{Замечание.}
В том же параграфе эта серия преобразуется к \textbf{интегралу Пуассона}:
\[
u(\rho,\varphi)
=
\frac{1}{2\pi}\int_0^{2\pi}
\frac{a^2-\rho^2}{a^2-2a\rho\cos(\varphi-\psi)+\rho^2}\,
g(\psi)\,d\psi.
\]
Для коллоквиума удобно помнить, что это не другой метод, а эквивалентная запись решения, уже найденного через ряд Фурье.

\paragraph{Мнемоническое правило / Суть на пальцах.}
Граничный профиль по углу сначала раскладывается на круговые гармоники \(\cos k\varphi\) и \(\sin k\varphi\). Каждая такая гармоника внутрь круга продолжается как \((\rho/a)^k\), то есть чем выше номер моды, тем быстрее её вклад затухает к центру.

\section*{Практические задачи к разделу 5}

\subsection*{Задача 5(a). Решить задачу Дирихле в круге радиуса \texorpdfstring{$R$}{R} при \texorpdfstring{$u(R,\varphi)=\cos\varphi$}{u(R,phi)=cos phi}}

\textbf{Дано.}
Найти функцию
\[
u(\rho,\varphi)
\]
из условий
\[
\Delta u=0
\qquad \text{в круге } 0\le \rho<R,
\]
\[
u(R,\varphi)=\cos\varphi.
\]

\textbf{Теоретическое обоснование.}
В пособии Алексеев использует обозначение $a$ для радиуса круга. В этой задаче нужно просто положить
\[
a=R.
\]
Тогда решение внутренней задачи Дирихле ищется по формуле
\[
u(\rho,\varphi)
=
\frac{\alpha_0}{2}
\sum_{k=1}^{\infty}
\left(\frac{\rho}{R}\right)^k
\bigl(\alpha_k\cos k\varphi+\beta_k\sin k\varphi\bigr),
\]
где $\alpha_k,\beta_k$ --- коэффициенты Фурье граничной функции
\[
g(\varphi)=\cos\varphi.
\]

\Solution
\textbf{Шаг 1. Разлагаем граничную функцию в ряд Фурье.}

В данном случае
\[
g(\varphi)=\cos\varphi.
\]
Это уже одна тригонометрическая гармоника. Поэтому её ряд Фурье имеет вид
\[
\cos\varphi=\frac{\alpha_0}{2}
\sum_{k=1}^{\infty}
\bigl(\alpha_k\cos k\varphi+\beta_k\sin k\varphi\bigr),
\]
где коэффициенты равны
\[
\alpha_1=1,
\qquad
\alpha_k=0 \text{ при } k\ne 1,
\qquad
\beta_k=0 \text{ для всех } k,
\qquad
\alpha_0=0.
\]

\textbf{Комментарий.} Здесь ничего интегрировать фактически не нужно, потому что граничная функция уже дана в базисной форме.

\textbf{Шаг 2. Подставляем коэффициенты в общую формулу решения.}

Так как ненулевым оказался только коэффициент при
\[
\cos\varphi,
\]
в ряде остаётся только гармоника с номером $k=1$:
\[
u(\rho,\varphi)
=
\left(\frac{\rho}{R}\right)\cos\varphi.
\]

\textbf{Шаг 3. Проверяем, что функция гармонична.}

Функция
\[
u(\rho,\varphi)=\frac{\rho}{R}\cos\varphi
\]
является одной из стандартных гармоник вида
\[
\rho^k\cos k\varphi
\]
при $k=1$, а значит, удовлетворяет уравнению Лапласа.

Если хочется увидеть это в декартовых координатах, то
\[
\rho\cos\varphi=x,
\]
поэтому
\[
u(x,y)=\frac{x}{R}.
\]
Тогда
\[
u_{xx}=0,
\qquad
u_{yy}=0,
\]
и, следовательно,
\[
\Delta u=u_{xx}+u_{yy}=0.
\]

\textbf{Шаг 4. Проверяем граничное условие.}

При $\rho=R$ имеем
\[
u(R,\varphi)=\frac{R}{R}\cos\varphi=\cos\varphi.
\]
Граничное условие выполнено.

\textbf{Ответ.}
\[
u(\rho,\varphi)=\frac{\rho}{R}\cos\varphi.
\]
Эквивалентно, в декартовых координатах:
\[
u(x,y)=\frac{x}{R}.
\]

\section*{6. Ряд Фурье и разложение по специальному базису}

\textit{Замечание по источнику. Этот раздел собран по нескольким местам пособия Г.\,В.~Алексеева.
Тригонометрические ряды на окружности используются в \S~6.4,
задача Штурма--Лиувилля и её свойства сформулированы в \S~4.2,
а примеры специальных базисов --- функций Лежандра и Бесселя ---
появляются в гл.~4 при разделении переменных в сферических и круговых областях.}

\subsection*{6(a). Определение тригонометрического ряда Фурье и формулы коэффициентов}

\Answer
\textbf{В пособии Алексеевa тригонометрический ряд Фурье естественно возникает на интервале}
\[
[0,2\pi]
\]
при разложении граничной функции $g(\varphi)$ на окружности. Именно в этой форме и запишем определение.

\textbf{Тригонометрическим рядом Фурье} функции $g$ называется ряд вида
\[
g(\varphi)\sim \frac{\alpha_0}{2}
+\sum_{k=1}^{\infty}\bigl(\alpha_k\cos k\varphi+\beta_k\sin k\varphi\bigr).
\]

\textbf{Коэффициенты Фурье} определяются формулами
\[
\alpha_0=\frac{1}{\pi}\int_0^{2\pi}g(\psi)\,d\psi,
\]
\[
\alpha_k=\frac{1}{\pi}\int_0^{2\pi}g(\psi)\cos k\psi\,d\psi,
\qquad
\beta_k=\frac{1}{\pi}\int_0^{2\pi}g(\psi)\sin k\psi\,d\psi,
\qquad
k=1,2,\dots
\]

\textbf{Почему именно такие формулы.}
Они получаются из ортогональности тригонометрической системы на $[0,2\pi]$:
\[
\int_0^{2\pi}\cos k\varphi \cos m\varphi\,d\varphi=0,
\qquad k\ne m,
\]
\[
\int_0^{2\pi}\sin k\varphi \sin m\varphi\,d\varphi=0,
\qquad k\ne m,
\]
\[
\int_0^{2\pi}\sin k\varphi \cos m\varphi\,d\varphi=0
\qquad \forall k,m.
\]

\textbf{Пошаговая схема вывода коэффициента \texorpdfstring{$\alpha_m$}{alpha m}.}

Предположим, что
\[
g(\varphi)=\frac{\alpha_0}{2}
+\sum_{k=1}^{\infty}\bigl(\alpha_k\cos k\varphi+\beta_k\sin k\varphi\bigr).
\]
Умножим обе части на $\cos m\varphi$ и проинтегрируем:
\[
\begin{aligned}
\int_0^{2\pi}g(\varphi)\cos m\varphi\,d\varphi
&=
\frac{\alpha_0}{2}\int_0^{2\pi}\cos m\varphi\,d\varphi
+\sum_{k=1}^{\infty}\alpha_k\int_0^{2\pi}\cos k\varphi\cos m\varphi\,d\varphi \\
&\quad
+\sum_{k=1}^{\infty}\beta_k\int_0^{2\pi}\sin k\varphi\cos m\varphi\,d\varphi.
\end{aligned}
\]
Первый интеграл равен нулю, все смешанные интегралы тоже равны нулю, а в сумме по косинусам остаётся только слагаемое при $k=m$:
\[
\int_0^{2\pi}g(\varphi)\cos m\varphi\,d\varphi
=
\alpha_m\int_0^{2\pi}\cos^2 m\varphi\,d\varphi.
\]
Поскольку
\[
\int_0^{2\pi}\cos^2 m\varphi\,d\varphi=\pi,
\]
получаем
\[
\alpha_m=\frac{1}{\pi}\int_0^{2\pi}g(\varphi)\cos m\varphi\,d\varphi.
\]
Аналогично выводится формула для $\beta_m$.

\textbf{Замечание об отрезке произвольной длины.}
Если функция задана на другом отрезке, то его обычно линейно переводят в $[0,2\pi]$ или используют специальные синусные и косинусные ряды на $(0,l)$:
\[
f(x)\sim \sum_{k=1}^{\infty}b_k\sin\frac{k\pi x}{l},
\qquad
b_k=\frac{2}{l}\int_0^l f(x)\sin\frac{k\pi x}{l}\,dx,
\]
или
\[
f(x)\sim \frac{a_0}{2}+\sum_{k=1}^{\infty}a_k\cos\frac{k\pi x}{l},
\qquad
a_k=\frac{2}{l}\int_0^l f(x)\cos\frac{k\pi x}{l}\,dx.
\]
Именно эти формы мы уже использовали в разделах 4 и 5.

\paragraph{Мнемоническое правило / Суть на пальцах.}
Ряд Фурье --- это способ разложить сложный профиль на сумму чистых гармоник: постоянной, косинусов и синусов. Коэффициенты просто измеряют, сколько каждой гармоники содержится в исходной функции.

\subsection*{6(b). Полнота и ортогональность. Примеры ортогональных систем}

\Answer
\textbf{Ортогональность.}
Пусть в пространстве функций задано скалярное произведение
\[
\langle f,g\rangle=\int_a^b w(x)f(x)g(x)\,dx,
\qquad w(x)>0.
\]
Система функций $\{\phi_k\}$ называется \textbf{ортогональной}, если
\[
\langle \phi_k,\phi_m\rangle=0
\qquad \text{при } k\ne m.
\]
Если дополнительно
\[
\langle \phi_k,\phi_k\rangle=1,
\]
то система называется \textbf{ортонормированной}.

\textbf{Полнота.}
Система $\{\phi_k\}$ называется \textbf{полной}, если с её помощью можно разложить достаточно широкий класс функций рассматриваемого пространства в ряд
\[
f\sim \sum_{k=1}^{\infty}c_k\phi_k,
\]
и если из условий
\[
\langle f,\phi_k\rangle=0
\qquad \forall k
\]
следует, что
\[
f\equiv 0
\]
в смысле выбранного пространства.

\textbf{Почему эти понятия важны.}
Ортогональность позволяет \textbf{вычислять коэффициенты разложения}, а полнота гарантирует, что система действительно достаточна для представления произвольной функции из нужного класса.

\textbf{Пример 1. Синусы на отрезке \texorpdfstring{$(0,l)$}{(0,l)}.}
Система
\[
\sin\frac{k\pi x}{l},
\qquad
k=1,2,\dots
\]
ортогональна на $(0,l)$:
\[
\int_0^l
\sin\frac{k\pi x}{l}\sin\frac{m\pi x}{l}\,dx=0
\qquad
(k\ne m).
\]
Именно эта система образует базис для задач с условиями Дирихле.

\textbf{Пример 2. Косинусы на отрезке \texorpdfstring{$(0,l)$}{(0,l)}.}
Система
\[
1,\quad \cos\frac{k\pi x}{l},
\qquad
k=1,2,\dots
\]
ортогональна на $(0,l)$ и естественно возникает при нулевых условиях Неймана.

\textbf{Пример 3. Полиномы Лежандра.}
В сферических задачах появляются полиномы Лежандра $P_n(x)$ и, более общо, присоединённые функции Лежандра $P_n^m(x)$. В книге показано, что они возникают лишь при
\[
\lambda_n=n(n+1),
\qquad
n=0,1,2,\dots
\]
и образуют ортогональную систему на $[-1,1]$. Для полиномов Лежандра это выглядит как
\[
\int_{-1}^1 P_n(x)P_m(x)\,dx=0
\qquad
(n\ne m).
\]

\textbf{Пример 4. Функции Бесселя.}
При разделении переменных в круге и цилиндрических областях возникает уравнение Бесселя. В задачах для круга появляются собственные функции вида
\[
J_\nu\!\left(\frac{\mu_{\nu,k}r}{R}\right),
\]
где $\mu_{\nu,k}$ --- корни соответствующей краевой задачи. Эти функции образуют ортогональную систему с весом $r$:
\[
\int_0^R
r\,
J_\nu\!\left(\frac{\mu_{\nu,k}r}{R}\right)
J_\nu\!\left(\frac{\mu_{\nu,m}r}{R}\right)\,dr=0
\qquad
(k\ne m).
\]

\textbf{Общий смысл.}
Во всех этих примерах один и тот же механизм:
\begin{itemize}
  \item область и граничные условия порождают спектральную задачу;
  \item её собственные функции оказываются ортогональными;
  \item полнота делает возможным разложение произвольных данных по этим функциям.
\end{itemize}

\paragraph{Мнемоническое правило / Суть на пальцах.}
Ортогональность --- это «гармоники не мешают друг другу». Полнота --- это «гармоник хватает, чтобы собрать любой допустимый профиль». Вместе эти свойства превращают систему функций в удобный рабочий базис.

\subsection*{6(c). Разложение по собственным функциям задачи Штурма--Лиувилля}

\Answer
\textbf{Общая задача Штурма--Лиувилля} в форме Алексеева возникает после разделения переменных и имеет вид
\[
[p(x)X'(x)]'+[\lambda \rho(x)-q(x)]X(x)=0
\qquad \text{на } (0,l),
\]
с граничными условиями
\[
\alpha X(0)-\beta X'(0)=0,
\qquad
\gamma X(l)+\delta X'(l)=0.
\]

\textbf{Собственные значения и собственные функции.}
Те значения параметра $\lambda$, при которых эта задача имеет нетривиальные решения, называются \textbf{собственными значениями}, а соответствующие нетривиальные решения $X_k$ --- \textbf{собственными функциями}.

\textbf{Шаг 1. Ортогональность.}

В книге доказывается, что если $\lambda_k\ne \lambda_m$, то собственные функции ортогональны с весом $\rho$:
\[
\int_0^l \rho(x)X_k(x)X_m(x)\,dx=0,
\qquad
k\ne m.
\]

\textbf{Шаг 2. Нормировка.}

Удобно либо оставить нормы в явном виде,
\[
\|X_k\|_\rho^2=\int_0^l \rho(x)X_k^2(x)\,dx,
\]
либо нормировать систему так, чтобы
\[
\int_0^l \rho(x)X_k^2(x)\,dx=1.
\]

\textbf{Шаг 3. Разложение функции.}

Если система собственных функций полна, то функцию $f(x)$ ищут в виде ряда
\[
f(x)\sim \sum_{k=1}^{\infty} c_k X_k(x).
\]
Чтобы найти $c_k$, умножаем разложение на $\rho(x)X_m(x)$ и интегрируем:
\[
\int_0^l \rho(x)f(x)X_m(x)\,dx
=
\sum_{k=1}^{\infty} c_k
\int_0^l \rho(x)X_k(x)X_m(x)\,dx.
\]
Благодаря ортогональности остаётся только одно слагаемое:
\[
\int_0^l \rho(x)f(x)X_m(x)\,dx
=
c_m\int_0^l \rho(x)X_m^2(x)\,dx.
\]
Отсюда
\[
c_m=
\frac{\displaystyle \int_0^l \rho(x)f(x)X_m(x)\,dx}
{\displaystyle \int_0^l \rho(x)X_m^2(x)\,dx}.
\]
Если система ортонормирована, формула упрощается:
\[
c_m=\int_0^l \rho(x)f(x)X_m(x)\,dx.
\]

\textbf{Как это связано с задачами математической физики.}
В уравнениях теплопроводности, колебаний и Лапласа неизвестная функция представляется как сумма мод:
\[
u(x,t)=\sum_{k=1}^{\infty}T_k(t)X_k(x).
\]
После этого:
\begin{itemize}
  \item пространственная часть $X_k$ определяется из спектральной задачи;
  \item временные коэффициенты $T_k$ находятся из начальных условий;
  \item сами начальные данные раскладываются в ряд по системе $\{X_k\}$.
\end{itemize}

\textbf{Итак, разложение по собственным функциям задачи Штурма--Лиувилля --- это прямой обобщённый аналог ряда Фурье.}
От обычного тригонометрического ряда оно отличается только тем, что базис строится не заранее, а порождается конкретной геометрией и конкретными краевыми условиями.

\paragraph{Мнемоническое правило / Суть на пальцах.}
Обычный ряд Фурье --- это частный случай. В общей задаче Штурма--Лиувилля сама физическая постановка «выбирает», какие гармоники разрешены, и именно по ним мы раскладываем данные.

\subsection*{6(d). Чем разложение в ряд Фурье отличается от разложения по специальным функциям}

\Answer
\textbf{Общее между ними.}
И ряд Фурье, и разложение по специальным функциям строятся по одной схеме:
\begin{itemize}
  \item выбирается ортогональная система функций;
  \item данные проектируются на эту систему;
  \item решение представляется как ряд по найденным коэффициентам.
\end{itemize}

\textbf{Главное различие --- в природе базиса.}

\textbf{Тригонометрический ряд Фурье} использует заранее фиксированную систему
\[
1,\ \cos kx,\ \sin kx
\]
или её варианты на отрезке $(0,l)$. Такой базис естественен для одномерных отрезков и круговой переменной.

\textbf{Разложение по специальным функциям} использует системы, которые появляются после разделения переменных в более сложной геометрии:
\begin{itemize}
  \item полиномы Лежандра --- в сферических и угловых задачах;
  \item функции Бесселя --- в круге, цилиндре, радиальных задачах;
  \item сферические функции --- в задачах на сфере.
\end{itemize}

\textbf{Различие в весе и в интервале ортогональности.}
У тригонометрических функций вес обычно равен $1$. У специальных функций появляются другие веса и другие интервалы:
\begin{itemize}
  \item для Лежандра интервалом служит $[-1,1]$;
  \item для Бесселя --- радиальный интервал $[0,R]$ с весом $r$;
  \item для общей задачи Штурма--Лиувилля весом служит функция $\rho(x)$.
\end{itemize}

\textbf{Различие в происхождении коэффициентов.}
Формально коэффициенты всегда находятся проекцией на базис, но сами базисные функции уже не обязаны быть синусами и косинусами. Они определяются тем дифференциальным уравнением, которое получилось после разделения переменных.

\textbf{Итоговая мысль.}
Тригонометрический ряд Фурье --- это специальный, самый простой случай разложения по собственным функциям. Специальные функции нужны там, где простая тригонометрическая система больше не согласована с геометрией области или с видом оператора.

\paragraph{Мнемоническое правило / Суть на пальцах.}
Синусы и косинусы --- это «базис для прямой и окружности». Лежандр, Бессель и другие специальные функции --- это «синусы и косинусы, приспособленные к новой геометрии».

\subsection*{6(e). Физический смысл коэффициентов Фурье в задачах теплопроводности и колебаний}

\Answer
\textbf{Коэффициенты Фурье} показывают, в какой мере исходное состояние содержит ту или иную собственную моду системы.

\textbf{В теплопроводности.}
Если
\[
u(x,t)=\sum_{k=1}^{\infty}a_k e^{-\omega_k^2 t}X_k(x),
\]
то коэффициент $a_k$ показывает, какой вклад мода $X_k$ вносит в начальное распределение температуры. Большие по модулю $a_k$ означают, что соответствующая пространственная гармоника сильно выражена в начальном профиле.

\textbf{Физический смысл затухания.}
Каждая мода живёт отдельно и затухает со своей скоростью. Поэтому коэффициенты Фурье описывают не просто форму начального состояния, но и то, какие компоненты исчезнут быстро, а какие будут доминировать при больших временах.

\textbf{В колебаниях.}
Для волновых задач решение имеет вид суммы стоячих волн:
\[
u(x,t)=\sum_{k=1}^{\infty}
\bigl(a_k\cos \omega_k t+b_k\sin \omega_k t\bigr)X_k(x).
\]
Здесь коэффициенты:
\begin{itemize}
  \item $a_k$ связаны с разложением начального смещения;
  \item $b_k$ связаны с разложением начальной скорости.
\end{itemize}

\textbf{Физически это означает:}
\begin{itemize}
  \item какие собственные частоты возбуждены;
  \item с какими амплитудами они возбуждены;
  \item какова роль каждой моды в общем движении.
\end{itemize}

\textbf{Общий вывод.}
Коэффициенты Фурье --- это координаты физического состояния в базисе собственных колебаний или собственных температурных мод.

\paragraph{Мнемоническое правило / Суть на пальцах.}
Коэффициенты Фурье отвечают на вопрос: «сколько каждой чистой моды сидит в данном процессе». В тепле это моды температуры, в колебаниях --- моды стоячих волн.

\section*{Практические задачи к разделу 6}

\subsection*{Задача 6(a). Разложить \texorpdfstring{$f(x)=x$}{f(x)=x} на \texorpdfstring{$(0,\pi)$}{(0,pi)} в ряд Фурье по синусам}

\textbf{Дано.}
Нужно разложить функцию
\[
f(x)=x,
\qquad
0<x<\pi,
\]
в синусный ряд Фурье
\[
x\sim \sum_{k=1}^{\infty} b_k \sin kx.
\]

\textbf{Теоретическое обоснование.}
На интервале $(0,\pi)$ синусный ряд имеет вид
\[
f(x)\sim \sum_{k=1}^{\infty}b_k\sin kx,
\qquad
b_k=\frac{2}{\pi}\int_0^{\pi}f(x)\sin kx\,dx.
\]
Эта форма естественна для задач с условиями Дирихле, потому что функции $\sin kx$ обращаются в нуль на концах $0$ и $\pi$.

\Solution
\textbf{Шаг 1. Записываем формулу коэффициентов.}

Подставляя $f(x)=x$, получаем
\[
b_k=\frac{2}{\pi}\int_0^{\pi}x\sin kx\,dx.
\]

\textbf{Шаг 2. Вычисляем интеграл по частям.}

Положим
\[
u=x,
\qquad
dv=\sin kx\,dx.
\]
Тогда
\[
du=dx,
\qquad
v=-\frac{\cos kx}{k}.
\]
По формуле интегрирования по частям
\[
\int_0^{\pi}x\sin kx\,dx
=
\left.-\frac{x\cos kx}{k}\right|_0^{\pi}
+\frac{1}{k}\int_0^{\pi}\cos kx\,dx.
\]

\textbf{Шаг 3. Вычисляем граничный член.}

Имеем
\[
\left.-\frac{x\cos kx}{k}\right|_0^{\pi}
=
-\frac{\pi\cos k\pi}{k}-0
=
-\frac{\pi(-1)^k}{k}.
\]

\textbf{Шаг 4. Вычисляем оставшийся интеграл.}

\[
\frac{1}{k}\int_0^{\pi}\cos kx\,dx
=
\frac{1}{k}\left.\frac{\sin kx}{k}\right|_0^{\pi}
=
\frac{1}{k^2}(\sin k\pi-\sin 0)=0.
\]

\textbf{Шаг 5. Собираем результат.}

Таким образом,
\[
\int_0^{\pi}x\sin kx\,dx
=
-\frac{\pi(-1)^k}{k}
=
\frac{\pi(-1)^{k+1}}{k}.
\]
Следовательно,
\[
b_k=\frac{2}{\pi}\cdot \frac{\pi(-1)^{k+1}}{k}
=
\frac{2(-1)^{k+1}}{k}.
\]

\textbf{Шаг 6. Записываем ряд.}

Итак,
\[
x=\sum_{k=1}^{\infty}\frac{2(-1)^{k+1}}{k}\sin kx,
\qquad
0<x<\pi.
\]

\textbf{Ответ.}
\[
x=2\sum_{k=1}^{\infty}\frac{(-1)^{k+1}}{k}\sin kx,
\qquad
0<x<\pi.
\]

\subsection*{Задача 6(b). Чем отличается разложение по синусам от разложения по косинусам и когда какое используется}

\textbf{Дано.}
Нужно сравнить два типа разложения функции на $(0,l)$:
\[
f(x)\sim \sum_{k=1}^{\infty}b_k\sin\frac{k\pi x}{l},
\]
и
\[
f(x)\sim \frac{a_0}{2}+\sum_{k=1}^{\infty}a_k\cos\frac{k\pi x}{l}.
\]

\textbf{Теоретическое обоснование.}
Синусы и косинусы --- это разные ортогональные системы, и выбор между ними определяется либо симметрией продолжения функции, либо типом граничных условий в задаче математической физики.

\Solution
\textbf{Шаг 1. Сравним сами базисные функции.}

Синусная система:
\[
\sin\frac{k\pi x}{l},
\qquad
k=1,2,\dots
\]
обращается в нуль на концах:
\[
\sin 0=0,
\qquad
\sin k\pi=0.
\]
Поэтому она естественна для задач с условиями
\[
u(0,t)=u(l,t)=0.
\]

Косинусная система:
\[
\cos\frac{k\pi x}{l},
\qquad
k=0,1,2,\dots
\]
имеет нулевую производную на концах:
\[
\frac{d}{dx}\cos\frac{k\pi x}{l}
=
-\frac{k\pi}{l}\sin\frac{k\pi x}{l},
\]
а значит,
\[
\left.\frac{d}{dx}\cos\frac{k\pi x}{l}\right|_{x=0}
=0,
\qquad
\left.\frac{d}{dx}\cos\frac{k\pi x}{l}\right|_{x=l}
=0.
\]
Поэтому косинусный ряд естественен для задач с условиями Неймана.

\textbf{Шаг 2. Сравним коэффициенты.}

Для синусного ряда
\[
b_k=\frac{2}{l}\int_0^l f(x)\sin\frac{k\pi x}{l}\,dx.
\]
Для косинусного ряда
\[
a_k=\frac{2}{l}\int_0^l f(x)\cos\frac{k\pi x}{l}\,dx,
\qquad
k\ge 1,
\]
\[
a_0=\frac{2}{l}\int_0^l f(x)\,dx.
\]

\textbf{Шаг 3. Свяжем это с продолжением функции.}

Синусный ряд соответствует \textbf{нечётному продолжению} функции с $(0,l)$ на $(-l,l)$.

Косинусный ряд соответствует \textbf{чётному продолжению} функции с $(0,l)$ на $(-l,l)$.

\textbf{Комментарий.} Именно поэтому в одном случае базис состоит только из синусов, а в другом --- только из косинусов.

\textbf{Шаг 4. Формулируем правило выбора.}

\begin{itemize}
  \item Если в задаче заданы условия Дирихле, особенно нулевые на концах, удобно использовать синусы.
  \item Если заданы условия Неймана, особенно нулевые производные на концах, удобно использовать косинусы.
  \item Если функция или постановка обладает нечётной симметрией, естествен синусный ряд.
  \item Если функция или постановка обладает чётной симметрией, естествен косинусный ряд.
\end{itemize}

\textbf{Ответ.}
Разложение по синусам и по косинусам отличается базисом, формулами коэффициентов и типом граничных условий, с которыми этот базис согласован. Синусы используют прежде всего для условий Дирихле и нечётного продолжения, косинусы --- для условий Неймана и чётного продолжения.

\subsection*{Задача 6(c). Задача Штурма--Лиувилля для \texorpdfstring{$d^2/dx^2$}{d2/dx2} с условиями Дирихле}

\textbf{Дано.}
Нужно рассмотреть задачу
\[
X''+\lambda X=0
\qquad \text{на } (0,l),
\]
\[
X(0)=0,
\qquad
X(l)=0.
\]

\textbf{Теоретическое обоснование.}
Это частный случай задачи Штурма--Лиувилля
\[
[p(x)X']'+[\lambda \rho(x)-q(x)]X=0
\]
при
\[
p(x)\equiv 1,
\qquad
\rho(x)\equiv 1,
\qquad
q(x)\equiv 0.
\]
Именно этот пример лежит в основе метода Фурье для одномерного теплового и волнового уравнений.

\Solution
\textbf{Шаг 1. Решаем уравнение при разных знаках \texorpdfstring{$\lambda$}{lambda}.}

\textbf{Случай \(\lambda=0\).}
Тогда
\[
X''=0,
\qquad
X(x)=C_1x+C_2.
\]
Из условий
\[
X(0)=0,
\qquad
X(l)=0
\]
следует
\[
C_2=0,
\qquad
C_1l=0,
\]
то есть
\[
X\equiv 0.
\]
Нетривиального решения нет.

\textbf{Случай \(\lambda<0\).}
Положим
\[
\lambda=-\mu^2,
\qquad
\mu>0.
\]
Тогда
\[
X''-\mu^2X=0,
\]
\[
X(x)=A\cosh\mu x+B\sinh\mu x.
\]
Из условия
\[
X(0)=0
\]
следует
\[
A=0.
\]
Тогда
\[
X(x)=B\sinh\mu x.
\]
Из условия
\[
X(l)=0
\]
получаем
\[
B\sinh\mu l=0.
\]
Поскольку
\[
\sinh\mu l\ne 0,
\]
остаётся только
\[
B=0.
\]
Нетривиального решения опять нет.

\textbf{Случай \(\lambda>0\).}
Положим
\[
\lambda=\mu^2,
\qquad
\mu>0.
\]
Тогда
\[
X''+\mu^2X=0,
\]
\[
X(x)=A\cos\mu x+B\sin\mu x.
\]
Из условия
\[
X(0)=0
\]
следует
\[
A=0.
\]
Значит,
\[
X(x)=B\sin\mu x.
\]
Теперь из
\[
X(l)=0
\]
имеем
\[
B\sin\mu l=0.
\]
Для нетривиального решения $B\ne 0$, поэтому
\[
\sin\mu l=0.
\]
Отсюда
\[
\mu l=k\pi,
\qquad
k=1,2,\dots
\]
и, следовательно,
\[
\mu_k=\frac{k\pi}{l},
\qquad
\lambda_k=\left(\frac{k\pi}{l}\right)^2.
\]
Соответствующие собственные функции:
\[
X_k(x)=\sin\frac{k\pi x}{l},
\qquad
k=1,2,\dots
\]

\textbf{Шаг 2. Проверяем ортогональность.}

Для $k\ne m$
\[
\int_0^l
\sin\frac{k\pi x}{l}\sin\frac{m\pi x}{l}\,dx=0.
\]
Таким образом, система собственных функций ортогональна в $L^2(0,l)$.

\textbf{Шаг 3. Нормировка.}

Кроме того,
\[
\int_0^l \sin^2\frac{k\pi x}{l}\,dx=\frac{l}{2}.
\]
Поэтому ортонормированная система имеет вид
\[
e_k(x)=\sqrt{\frac{2}{l}}\sin\frac{k\pi x}{l},
\qquad
k=1,2,\dots
\]

\textbf{Шаг 4. Почему это полный ортогональный базис.}

Это регулярная задача Штурма--Лиувилля с весом
\[
\rho(x)\equiv 1.
\]
Общая теория, использованная в книге при методе Фурье, утверждает, что все собственные функции такой задачи образуют полную ортогональную систему. А в нашем частном случае это означает, что любую достаточно хорошую функцию на $(0,l)$, согласованную с условиями Дирихле, можно разложить в синусный ряд
\[
f(x)\sim \sum_{k=1}^{\infty}c_k\sin\frac{k\pi x}{l}.
\]

\textbf{Ответ.}
Собственные значения:
\[
\lambda_k=\left(\frac{k\pi}{l}\right)^2,
\qquad
k=1,2,\dots
\]
Собственные функции:
\[
X_k(x)=\sin\frac{k\pi x}{l},
\qquad
k=1,2,\dots
\]
Они образуют полный ортогональный базис в $L^2(0,l)$, поскольку это собственные функции регулярной задачи Штурма--Лиувилля с весом $\rho\equiv 1$.

\section*{7. Сравнение трёх типов уравнений}

\textit{Замечание по источнику. Этот раздел является итоговой синтезирующей сводкой по гл.~2--6 пособия Алексеева: классификация на эллиптические, параболические и гиперболические уравнения обсуждается в гл.~2, а характерные свойства затем подробно иллюстрируются на уравнениях Пуассона, теплопроводности и волновом уравнении.}

\subsection*{7(a). Основные свойства решений эллиптических, параболических и гиперболических уравнений}

\Answer
\textbf{Простейшие представители трёх типов.}
\begin{itemize}
  \item \textbf{Эллиптический тип:} уравнение Лапласа или Пуассона.
  \item \textbf{Параболический тип:} уравнение теплопроводности.
  \item \textbf{Гиперболический тип:} волновое уравнение.
\end{itemize}

\textbf{1. Скорость распространения возмущений.}

\textbf{Эллиптические уравнения.}
Они описывают \textbf{стационарные} состояния. Поэтому вопрос о скорости распространения во времени здесь не является основным. Изменение граничных данных меняет решение сразу во всей области, поскольку задача ставится на весь пространственный домен целиком.

\textbf{Параболические уравнения.}
Для уравнения теплопроводности характерна \textbf{бесконечная скорость распространения возмущений}. Даже локальное начальное возмущение при любом $t>0$ влияет на решение во всей области.

\textbf{Гиперболические уравнения.}
Для волнового уравнения характерна \textbf{конечная скорость распространения}. Возмущение распространяется по характеристикам и фронтам волн, а не мгновенно по всей области.

\textbf{2. Гладкость решений.}

\textbf{Эллиптические уравнения.}
Гармонические функции обладают высокой внутренней гладкостью; внутри области решение значительно более гладко, чем исходные данные на границе.

\textbf{Параболические уравнения.}
Решение уравнения теплопроводности при $t>0$ мгновенно сглаживается: высокочастотные моды быстро затухают, поэтому решение становится очень гладким.

\textbf{Гиперболические уравнения.}
Сглаживающий эффект, как правило, отсутствует. Негладкость исходных данных не исчезает мгновенно, а переносится и распространяется вместе с волновым процессом.

\textbf{3. Принцип максимума.}

\textbf{Эллиптические уравнения.}
Для гармонических функций справедлив принцип максимума: максимум и минимум достигаются на границе области.

\textbf{Параболические уравнения.}
Для уравнения теплопроводности также действует принцип максимума, но уже в пространственно-временном цилиндре: экстремум контролируется начальными и граничными данными.

\textbf{Гиперболические уравнения.}
Аналога принципа максимума в таком виде нет. Поведение решения определяется не только граничными значениями, но и начальной скоростью, а также законом распространения по характеристикам.

\textbf{4. Наличие характеристик.}

\textbf{Эллиптические уравнения.}
В классическом смысле реальных характеристик, играющих роль линий распространения информации, нет.

\textbf{Параболические уравнения.}
Характеристики в формальном смысле вырождаются и не играют той роли, какую они играют у гиперболических уравнений. Основную роль здесь играют принцип максимума и сглаживание.

\textbf{Гиперболические уравнения.}
Характеристики являются центральным объектом: именно вдоль них переносится информация и строится анализ решения.

\textbf{5. Физическая интерпретация.}

\begin{itemize}
  \item Эллиптические уравнения описывают \textbf{установившиеся поля}.
  \item Параболические уравнения описывают \textbf{диффузию и выравнивание}.
  \item Гиперболические уравнения описывают \textbf{волны и инерционные процессы}.
\end{itemize}

\paragraph{Мнемоническое правило / Суть на пальцах.}
Эллиптическое уравнение отвечает на вопрос «какой профиль установился». Параболическое --- «как всё растекается и сглаживается». Гиперболическое --- «как всё бежит и колеблется с конечной скоростью».

\subsection*{7(b). Какие типы задач ставятся для каждого типа}

\Answer
\textbf{Для эллиптических уравнений} обычно ставят \textbf{краевые задачи}, потому что время в них не участвует как эволюционный параметр. Основные примеры:
\begin{itemize}
  \item задача Дирихле;
  \item задача Неймана;
  \item третья краевая задача;
  \item смешанная краевая задача.
\end{itemize}

\textbf{Для параболических уравнений} обычно ставят \textbf{начально-краевые задачи}, поскольку требуется:
\begin{itemize}
  \item одно начальное условие по времени;
  \item граничные условия по пространственным переменным.
\end{itemize}
На всей прямой или в неограниченной области может ставиться и задача Коши.

\textbf{Для гиперболических уравнений} ставят:
\begin{itemize}
  \item задачу Коши;
  \item смешанные начально-краевые задачи;
  \item задачи на характеристиках.
\end{itemize}
Из-за второго порядка по времени здесь, как правило, требуется \textbf{два начальных условия}: для самого решения и для его производной по времени.

\textbf{Сводка по числу исходных данных.}
\begin{itemize}
  \item Эллиптическое уравнение: задаются граничные данные на всей границе.
  \item Параболическое уравнение: задаются одно начальное условие и граничные данные.
  \item Гиперболическое уравнение: задаются два начальных условия и, при необходимости, дополнительные граничные условия.
\end{itemize}

\paragraph{Мнемоническое правило / Суть на пальцах.}
Для эллиптики важна граница, для параболики --- старт плюс граница, для гиперболики --- стартовое положение, стартовая скорость и дальше уже геометрия распространения волны.

\end{document}
